% !TeX root = ../../Book4.tex
% 纲要标题：## 第9章　Hamilton 动力学、通量与 Poisson 结构
% 纲要位置：几何学计划纲要.md，第 3812 行。
% 正文状态：初稿；尚未正式审核。

% 纲要内容开始
% 
% > 本章目标与证明范围：由作用变分进入 Hamilton 系统，再用闭与恰当一形式的区别研究辛同位。Poisson 的退化情况是一条独立延伸，不与非退化辛情形混写。
% 
% 纲要内容结束

\chapter{Hamilton 动力学、通量与 Poisson 结构}
\label{g4:ch:09}
\BookChapterNumber{g4:ch:09}

\begin{chapterabstract}
辛形式把函数的微分转成向量场，使能量函数产生运动。本章固定 Hamilton 向量场的正号约定，从余切坐标推导运动方程，再通过作用量的一阶变分说明这些方程如何依赖端点条件。Poisson 括号记录观测量沿运动的变化，Jacobi 恒等式把守恒量、向量场括号与对称联系起来。闭一形式和恰当一形式的区别进一步给出辛同位的通量，并在环面上形成可直接计算的整体障碍。最后允许 Poisson 张量退化，证明局部分裂及辛叶的存在，用线性 Lie--Poisson 结构和秩跳跃例子展示不同维数的动力学区域。
\end{chapterabstract}

\begin{learninggoals}
\begin{itemize}
	\item 从 \(\iota_{X_H}\omega=\mathrm{d}H\) 推导 Hamilton 方程，计算基本轨道，并区分能量守恒与流的完备性。
	\item 完整计算作用量的一阶变分，分别判断固定相点、固定底点及自由端点给出的条件。
	\item 证明 Poisson 括号的 Jacobi 恒等式，核对它与向量场括号之间的负号，并用括号寻找第一积分及检验约束。
	\item 通过闭与恰当一形式判别辛场和 Hamilton 场，构造光滑生成函数族，理解路径与终点的区别。
	\item 定义并计算通量，证明端点固定同伦下的不变性，求出标准环面的通量群与平移障碍。
	\item 用坐标检验 Poisson 条件，证明局部分裂及辛叶结构，计算线性模型、奇异叶和 Casimir 函数。
\end{itemize}
\end{learninggoals}

一个粒子的位置满足二阶运动方程。加入动量后，同样的运动可以写成相空间中的一阶方程：位置变化由能量对动量的导数给出，动量变化由能量对位置的负导数给出。这个看似依赖坐标的配对来自余切丛的辛形式。因此，Hamilton 系统可以从具体的力学方程推广到任意辛流形上的光滑函数。

\ProseParagraphBreak

函数在这里承担两种作用：它可以生成运动，也可以用来观测另一种运动。将前者产生的向量场作用到后者，就得到 Poisson 括号。一个观测量的括号为零，意味着它沿相应运动不变；若有多个独立的对合积分，轨道便被限制在较小的水平集中。向量场是否交换、约束是否保持，也都能由同一个括号计算。

\ProseParagraphBreak

局部方程还留下一个整体问题。环面上的平移处处保持面积，每个小坐标图内都能找到线性 Hamilton 函数，但这个函数可能绕一圈后改变一个常数。闭一形式的周期记录了这个障碍。沿一整条辛同位累积这些周期，就得到通量；改变连接同一终点的路径，还可能加上一份环路通量。把两个层次分开，才能判断一个终点是否存在任何 Hamilton 生成路径。

\ProseParagraphBreak

本章使用第七章的余切丛与辛线性代数、第八章的 Darboux 坐标及随流拉回公式，以及第一卷的光滑流和正则水平集。上同调只在通量部分用于闭一形式模恰当形式；环面计算用到第二卷已得的一次上同调。\secref{g4:sec:0904}从函数括号出发允许秩下降，并在本节内证明所需局部模型与辛叶结论。Hamilton 场、Poisson 括号、余伴随轨道及第一积分的正式参考见 Cannas da Silva 第5.1、5.5节；涉及动力学时始终保留解的存在区间。\cite[第5.1、5.5节]{Cannas2006Symplectic}

% !TeX root = ../../Book4.tex
% 纲要标题：### 9.1 Hamilton 方程与作用
% 纲要位置：几何学计划纲要.md，第 3816 行。
% 正文状态：初稿；尚未正式审核。

% 纲要内容开始
% 
% * ι_XH ω=dH 的固定约定
% * 余切坐标中的 Hamilton 方程
% * 作用及端点条件
% * 时间依赖 Hamilton 和守恒量
% 
% 纲要内容结束

\section{Hamilton 方程与作用}
\label{g4:sec:0901}

位置与动量一起描述运动状态，辛形式则把一个函数的微分转成状态空间中的速度。这个转换给出 Hamilton 方程。另一方面，余切丛的典范一形式可以沿轨迹积分，得到比较不同轨迹的作用量。本节将证明，两种描述给出同一组运动方程；符号与端点条件会在变分计算中同时确定。

\subsection{\texorpdfstring{\(\iota_{X_H}\omega=\mathrm{d}H\) 的固定约定}{\GeometryPdfIota{}\_XH \GeometryPdfSmallOmega{}=dH 的固定约定}}
\label{g4:subsec:0901:01}

在辛流形 \((M,\omega)\) 上，每点的映射 \(v\mapsto\omega(v,\cdot)\) 是切空间到余切空间的同构。因此一个光滑函数 \(H\) 的微分，可以唯一决定一个向量场。这里插入始终放在二形式的第一个变量。

\begin{definition}\label{g4:def:0901:hamiltonian-field}
	设 \((M,\omega)\) 为无边界辛流形，\(H\in C^\infty(M,\mathbb R)\)。由
	\begin{equation}\label{g4:eq:0901:hamiltonian-convention}
		\iota_{X_H}\omega=\mathrm{d}H,
		\qquad \omega(X_H,v)=\mathrm{d}H(v)
	\end{equation}
	确定的向量场 \(X_H\) 称为 \(H\) 的\term[hamiltonxiangliangchang]{Hamilton 向量场}[Hamiltonian vector field]\footnote{William Rowan Hamilton（1805--1865），爱尔兰数学家，研究力学与光学，并创立四元数理论。}，\(H\) 称为该场的\term[hamiltonhanshu]{Hamilton 函数}[Hamiltonian function]。积分曲线 \(\dot z=X_H(z)\) 称为 Hamilton 轨迹；它们组成的局部流称为\term[hamiltonliu]{Hamilton 流}[Hamiltonian flow]。
\end{definition}

函数 \(H\) 不一定来自某个具体力学系统；能量函数是其中一个重要来源。定义本身只使用辛形式和函数的微分，无需另选度量。下面先核对它的光滑性和最基本的性质。

\begin{proposition}\label{g4:prop:0901:basic-properties}
	设 \((M,\omega)\) 为无边界辛流形。每个光滑实函数 \(H\) 都唯一确定一个光滑 Hamilton 向量场，且对光滑实函数 \(H,K\) 及 \(a,b\in\mathbb R\) 有
	\[
		X_{aH+bK}=aX_H+bX_K.
	\]
	此外，\(X_H(x)=0\) 当且仅当 \(\mathrm{d}H_x=0\)；\(X_H=X_K\) 当且仅当 \(H-K\) 在每个连通分支上为常数。这个场满足
	\begin{equation}\label{g4:eq:0901:hamiltonian-invariance}
		\mathrm{d}H(X_H)=0,
		\qquad \mathcal L_{X_H}\omega=0.
	\end{equation}
\end{proposition}
\begin{proof}
	逐点的非退化性给出唯一解 \(X_H=(\omega^\flat)^{-1}\mathrm{d}H\)。在局部标架中，\(\omega^\flat\) 是光滑可逆矩阵，其逆由伴随矩阵除以非零行列式表示，所以所得场光滑。线性性及零点判别也由这个同构得到。\(X_H=X_K\) 等价于 \(\mathrm{d}(H-K)=0\)，而微分为零的光滑函数在每个连通分支上为常数。

	反对称性给出 \(\mathrm{d}H(X_H)=\omega(X_H,X_H)=0\)。由第一卷定理14.1.5的 Cartan 公式，以及 \(\mathrm{d}\omega=0\)，
	\[
		\mathcal L_{X_H}\omega
		=\mathrm{d}(\iota_{X_H}\omega)+\iota_{X_H}\mathrm{d}\omega
		=\mathrm{d}^2H=0.
	\]
	这证明所列性质。
\end{proof}

第一式说明沿自主 Hamilton 轨迹函数 \(H\) 不变；第二式是辛形式在无穷小意义下不变。它们使用的条件有所不同：前者来自反对称性，后者还用到辛形式的闭性。把无穷小不变性积分成流的不变性时，需要同时保留流的定义域，末小节将明确说明。上述定义和基本性质可对照 Cannas da Silva 的第5.1节。\cite[第139--140页]{Cannas2006Symplectic}

\subsection{余切坐标中的 Hamilton 方程}
\label{g4:subsec:0901:02}

在余切丛 \(\mathrm{T}^*Q\) 上，\cref{g4:thm:0702:cotangent-symplectic}给出
\[
	\theta=\sum_{i=1}^n p_i\,\mathrm{d}q^i,
	\qquad \omega=-\mathrm{d}\theta
	=\sum_{i=1}^n\mathrm{d}q^i\wedge\mathrm{d}p_i.
\]
称 \(q^i\) 为位置坐标，\(p_i\) 为与之配对的动量坐标。力学中把同时记录这些变量的空间称为\term[xiangkongjian]{相空间}[phase space]；底流形 \(Q\) 则记录位置或构形。

\begin{proposition}\label{g4:prop:0901:hamiltonian-coordinates}
	设 \(H\) 是余切丛 \(\mathrm{T}^*Q\) 上的光滑实函数。在任一余切坐标域内，
	\begin{equation}\label{g4:eq:0901:hamiltonian-field-coordinates}
		X_H=\sum_{i=1}^n\left(
		\frac{\partial H}{\partial p_i}\frac{\partial}{\partial q^i}
		-\frac{\partial H}{\partial q^i}\frac{\partial}{\partial p_i}
		\right).
	\end{equation}
	因此曲线 \(z(t)=(q(t),p(t))\) 是 Hamilton 轨迹，当且仅当它满足
	\begin{equation}\label{g4:eq:0901:hamilton-equations}
		\dot q^i=\frac{\partial H}{\partial p_i},
		\qquad \dot p_i=-\frac{\partial H}{\partial q^i},
		\qquad 1\le i\le n.
	\end{equation}
	同一表达适用于任意辛流形上的 Darboux 坐标。
\end{proposition}
\begin{proof}
	写 \(X_H=\sum_i(a^i\partial_{q^i}+b_i\partial_{p_i})\)，则
	\[
		\iota_{X_H}\omega=\sum_i(a^i\,\mathrm{d}p_i-b_i\,\mathrm{d}q^i).
	\]
	与 \(\mathrm{d}H=\sum_i(H_{q^i}\,\mathrm{d}q^i+H_{p_i}\,\mathrm{d}p_i)\) 比较系数，得 \(a^i=H_{p_i}\)、\(b_i=-H_{q^i}\)。代入积分曲线方程即得结论。计算只用了辛形式的标准表达，所以同样适用于 Darboux 坐标。
\end{proof}

在标准 Euclid 内积下，这个速度场通常不指向 \(H\) 增长最快的方向。按 \((q,p)\) 顺序记 \(J_0(q,p)=(-p,q)\)，则式\eqref{g4:eq:0901:hamiltonian-field-coordinates}给出 \(X_H=-J_0\nabla H\)。在平面上，它把梯度旋转四分之一周后沿等值线运动；方向由式\eqref{g4:eq:0901:hamiltonian-convention}固定。

\begin{example}\label{g4:exm:0901:mechanical-system}
	在 \(\mathrm{T}^*\mathbb R^n\) 上取
	\[
		H(q,p)=\frac{|p|^2}{2m}+V(q),\qquad m>0,
	\]
	其中 \(V\) 为光滑势能。Hamilton 方程为
	\[
		\dot q=\frac pm,\qquad \dot p=-\nabla V(q).
	\]
	消去 \(p=m\dot q\)，得到 \(m\ddot q=-\nabla V(q)\)。反过来，任一满足这条二阶方程的 \(q\)，连同 \(p=m\dot q\)，都给出 Hamilton 轨迹。因此在位置空间中的二阶运动，可以转成相空间中的一阶运动。质量 \(m\) 把速度转为动量，而 \(H\) 沿运动保持为动能与势能之和。\cite[第5.1节，第140页]{Cannas2006Symplectic}
\end{example}

\begin{example}\label{g4:exm:0901:harmonic-oscillator}
	在 \(\mathbb R^2_{q,p}\) 中取 \(H=p^2/(2m)+kq^2/2\)，其中 \(m,k>0\)，并令 \(\nu=\sqrt{k/m}\)。方程 \(\dot q=p/m\)、\(\dot p=-kq\) 给出
	\[
		\begin{aligned}
			q(t)&=q_0\cos(\nu t)+\frac{p_0}{m\nu}\sin(\nu t),\\
			p(t)&=p_0\cos(\nu t)-m\nu q_0\sin(\nu t).
		\end{aligned}
	\]
	求导及代入 \(t=0\) 可逐项验证方程与初值，唯一性因而说明这就是全体轨迹。所有时间都有定义；原点是静止点，正能量等值线是椭圆，非静止轨迹的周期为 \(2\pi/\nu\)。在 \(m=k=1\) 时，从 \((1,0)\) 出发的初速度是 \((0,-1)\)，所以相平面中的运动方向为顺时针。这与 \(X_H=-J_0\nabla H\) 一致。
\end{example}

\subsection{作用及端点条件}
\label{g4:subsec:0901:03}

固定两个时刻 \(a<b\)。对 \(\mathrm{T}^*Q\) 中的光滑曲线 \(z(t)=(q(t),p(t))\)，典范一形式给出 \(\theta(\dot z)=p(\dot q)\)。若允许 \(H=H(t,z)\) 随时间变化，定义其\term[zuoyongliang]{作用量}[action functional]为
\begin{equation}\label{g4:eq:0901:action}
	\mathcal A_H(z)=\int_a^b\bigl(\theta_{z(t)}(\dot z(t))-H(t,z(t))\bigr)\,\mathrm{d}t.
\end{equation}
在余切坐标中，这就是 \(\int_a^b(\sum_i p_i\dot q^i-H)\,\mathrm{d}t\)。由于 \(\theta\) 在整个余切丛上定义，曲线不必留在一张坐标图中。下面的计算也适用于任一已经给定一形式 \(\theta\) 且 \(\omega=-\mathrm{d}\theta\) 的辛流形。一般辛流形未必有这样的整体 \(\theta\)，此时不能直接把式\eqref{g4:eq:0901:action}作为所有曲线上的整体定义。

\ProseParagraphBreak

一个光滑变分是映射 \(Z:(-\varepsilon,\varepsilon)\times[a,b]\to M\)，记 \(z_s(t)=Z(s,t)\)，并要求 \(z_0=z\)。沿 \(z\) 的变分向量场为 \(V(t)=\partial_sZ(0,t)\)。曲线相对于一类允许变分为\term[zuoyonglianglinjiequxian]{作用量的临界曲线}[critical curve of the action]，是指所有这些变分都满足 \(\left.\mathrm{d}\mathcal A_H(z_s)/\mathrm{d}s\right|_{s=0}=0\)。端点条件规定的正是哪些变分允许使用。本小节中的时刻 \(a,b\) 始终固定。

\begin{theorem}\label{g4:thm:0901:first-variation}
	设 \((M,\omega)\) 为无边界辛流形，且给定 \(\theta\in\Omega^1(M)\) 满足 \(\omega=-\mathrm{d}\theta\)。设 \(H:[a,b]\times M\to\mathbb R\) 光滑，\(X_{H_t}\) 按空间微分 \(\iota_{X_{H_t}}\omega=\mathrm{d}_M H_t\) 定义。对任意光滑曲线 \(z\) 及光滑变分 \(z_s\)，有
	\begin{equation}\label{g4:eq:0901:first-variation}
		\left.\frac{\mathrm{d}}{\mathrm{d}s}\right|_{s=0}\mathcal A_H(z_s)
		=\bigl[\theta(V)\bigr]_a^b
		+\int_a^b\omega\bigl(\dot z-X_{H_t}(z),V\bigr)\,\mathrm{d}t.
	\end{equation}
	这里 \([\theta(V)]_a^b=\theta_{z(b)}(V(b))-\theta_{z(a)}(V(a))\)。在余切坐标中，积分项等于
	\begin{equation}\label{g4:eq:0901:first-variation-coordinates}
		\int_a^b\sum_i\left[
		\left(\dot q^i-\frac{\partial H}{\partial p_i}\right)\delta p_i
		-\left(\dot p_i+\frac{\partial H}{\partial q^i}\right)\delta q^i
		\right]\mathrm{d}t,
	\end{equation}
	边界项为 \([\sum_i p_i\delta q^i]_a^b\)。
\end{theorem}
\begin{proof}
	在参数矩形上写
	\[
		Z^*\theta=A(s,t)\,\mathrm{d}t+B(s,t)\,\mathrm{d}s,
	\]
	其中 \(A=\theta(\partial_tZ)\)、\(B=\theta(\partial_sZ)\)。第一卷命题14.1.8的拉回与外微分相容性给出
	\[
		\partial_sA-\partial_tB
		=(Z^*\mathrm{d}\theta)(\partial_s,\partial_t)
		=\mathrm{d}\theta(\partial_sZ,\partial_tZ).
	\]
	因此，先对 \(s\) 求导再沿 \(t\) 积分，得到
	\[
		\left.\frac{\mathrm{d}}{\mathrm{d}s}\right|_0
		\int_a^b\theta(\dot z_s)\,\mathrm{d}t
		=\bigl[\theta(V)\bigr]_a^b
		+\int_a^b\mathrm{d}\theta(V,\dot z)\,\mathrm{d}t.
	\]
	被积函数联合光滑且时间区间紧，所以这里可以把微分移入积分。Hamilton 项的导数为 \(-\int_a^b\mathrm{d}_M H_t(V)\,\mathrm{d}t\)。使用 \(\mathrm{d}\theta=-\omega\) 与 \(\mathrm{d}_M H_t(V)=\omega(X_{H_t},V)\)，便有
	\[
		\mathrm{d}\theta(V,\dot z)-\mathrm{d}_M H_t(V)
		=\omega(\dot z,V)-\omega(X_{H_t},V),
	\]
	得到式\eqref{g4:eq:0901:first-variation}。

	在余切坐标中，\(V=\sum_i(\delta q^i\partial_{q^i}+\delta p_i\partial_{p_i})\)。代入 \(\omega=\sum_i\mathrm{d}q^i\wedge\mathrm{d}p_i\) 与式\eqref{g4:eq:0901:hamiltonian-field-coordinates}，就得到式\eqref{g4:eq:0901:first-variation-coordinates}；典范一形式的表达给出所列边界项。也可以直接对 \(\sum_i p_i\dot q^i-H\) 求变分，再将 \(p_i\,\delta\dot q^i\) 分部积分，获得相同结果。
\end{proof}

\begin{proposition}\label{g4:prop:0901:endpoint-conditions}
	在\cref{g4:thm:0901:first-variation}的条件下，设 \(P_a,P_b\subset M\) 为嵌入子流形，\(z(a)\in P_a\)、\(z(b)\in P_b\)，并只考虑满足 \(z_s(a)\in P_a\)、\(z_s(b)\in P_b\) 的变分。则 \(z\) 为作用量的临界曲线，当且仅当
	\[
		\dot z=X_{H_t}(z),\qquad
		\theta_{z(a)}|_{\mathrm{T}_{z(a)}P_a}=0,\qquad
		\theta_{z(b)}|_{\mathrm{T}_{z(b)}P_b}=0.
	\]
\end{proposition}
\begin{proof}
	先取支集紧含于 \((a,b)\) 的局部坐标变分，它不改变端点，因此始终允许。式\eqref{g4:eq:0901:first-variation}的边界项为零，故一形式 \(\omega(\dot z-X_{H_t},\cdot)\) 沿曲线与所有这样的测试向量场积分配对为零。若其中一个坐标分量在某时刻非零，缩小到其保持同号的时间段，并乘一个非负非零光滑截断函数作为该坐标方向的变分，就会得到非零积分，矛盾。因此这个一形式处处为零，非退化性给出 \(\dot z=X_{H_t}(z)\)。这些测试向量场可在坐标中通过 \(z_s=z+sV\) 实现，支集外保持原曲线，故确实来自允许的光滑变分。

	沿 Hamilton 轨迹，变分只剩两端的边界项。端点在嵌入子流形内运动时，其一阶速度可以是该子流形的任意切向量：在适配坐标中选取具有指定初速度的端点曲线，再用时间截断函数把这一改动延入靠近端点的一小段，即得相应变分。两端可独立变化，所以驻值条件恰好要求两项分别消失。反过来，运动方程使积分项为零，所列端点条件使边界项为零，因此任意允许变分的一阶作用变化都为零。
\end{proof}

在余切丛中，这个统一结论给出三种常用边界条件。
\begin{enumerate}
	\item 若固定两个相空间端点 \(z(a),z(b)\)，则 \(V(a)=V(b)=0\)，边界项自动消失。因此临界曲线恰为具有这些端点的 Hamilton 轨迹。
	\item 若只固定底空间端点 \(q(a),q(b)\)，动量端点仍可变化，但 \(\delta q(a)=\delta q(b)=0\)，所以 \([p\,\delta q]_a^b=0\) 仍自动成立。几何上，端点位于两个固定余切纤维，\(\theta\) 在竖直切向量上本来就为零。此时同样只得到 Hamilton 方程，没有额外的端点动量条件。
	\item 若两个相空间端点都完全自由，则允许任意端点切向量，必须有 \(\theta_{z(a)}=\theta_{z(b)}=0\)。由典范一形式的定义，这等价于 \(p(a)=p(b)=0\)。这些自然端点条件与 Hamilton 方程一起构成临界曲线的条件。
\end{enumerate}

更一般地，若只要求底端点落在嵌入子流形 \(C_a,C_b\subset Q\) 中，而端点动量不受限制，则自然条件为 \(p(a)|_{\mathrm{T}_{q(a)}C_a}=0\)、\(p(b)|_{\mathrm{T}_{q(b)}C_b}=0\)。固定底端点与完全自由底端点分别对应 \(C\) 为一点和整个 \(Q\)。这里的自由始终指固定时刻的空间端点可以移动。

\ProseParagraphBreak

变分等价本身不保证任意指定的两端点都能由 Hamilton 轨迹连接。一阶微分方程的初值已经确定整条轨迹，再任意指定另一相空间端点可能无解。例如 \(H=0\) 时所有轨迹恒定，两个不同相空间端点之间没有临界曲线。把两端只固定在底空间，保留了选择初始动量的自由，因而是力学边值问题常用的条件。

\begin{example}\label{g4:exm:0901:stationary-not-minimum}
	取 \(M=\mathbb R^2\)、\(H=0\)、\([a,b]=[0,1]\)，并固定两个相空间端点为原点。常值曲线 \(z=0\) 是 Hamilton 轨迹，所以是作用量的临界曲线。对变分
	\[
		q_s(t)=s\sin(\pi t),\qquad
		p_s(t)=\pm s\sin(2\pi t),
	\]
	端点始终为原点，但
	\[
		\mathcal A_0(z_s)
		=\pm s^2\pi\int_0^1\sin(2\pi t)\cos(\pi t)\,\mathrm{d}t
		=\pm\frac43s^2.
	\]
	这两族变分可以任意小，却分别使作用量变大和变小。因此这里得到的是驻值原理，不能把每条 Hamilton 轨迹都称为作用量的极小曲线。
\end{example}

对于\cref{g4:exm:0901:mechanical-system}，先对动量作变分得到 \(p=m\dot q\)，代回作用量可得
\[
	\int_a^b\left(\frac m2|\dot q|^2-V(q)\right)\,\mathrm{d}t.
\]
在固定底端点的变分下，对动能项分部积分，所得一阶变分为 \(-\int_a^b(m\ddot q+\nabla V)\cdot\delta q\,\mathrm{d}t\)，再次得到 Newton 方程。\footnote{Isaac Newton（1643--1727），英国数学家与物理学家，研究微积分、力学与光学。}相空间的一阶方程与位置空间的二阶方程因而具有相同的变分来源。固定端点的分部积分论证可对照 Cannas da Silva 的第5.3节；其第5.2节的作用微分也采用 \(\omega(\dot z-X_H,V)\) 的次序。\cite[第142--146页]{Cannas2006Symplectic}

\subsection{时间依赖 Hamilton 和守恒量}
\label{g4:subsec:0901:04}

现在在一般辛流形 \((M,\omega)\) 上固定一个光滑函数 \(H:I\times M\to\mathbb R\)，其中 \(I\subset\mathbb R\) 是开时间区间。记 \(H_t(x)=H(t,x)\)，逐时刻定义
\[
	\iota_{X_{H_t}}\omega=\mathrm{d}_M H_t.
\]
右端只对空间求微分，不包含 \((\partial_tH)\,\mathrm{d}t\)。矩阵求逆的光滑性保证 \(X_{H_t}\) 对 \((t,x)\) 联合光滑。若 \(H\) 不显含时间，称为\term[zizhuHamiltonxitong]{自主 Hamilton 系统}[autonomous Hamiltonian system]；一般情形称为时间依赖 Hamilton 系统。

\ProseParagraphBreak

由第一卷命题13.1.9，初值 \(z(s)=x\) 确定唯一最大解，记其在时刻 \(t\) 的值为 \(\Phi_{t,s}(x)\)。演化映射在开放的定义域上光滑，满足 \(\Phi_{t,r}\Phi_{r,s}=\Phi_{t,s}\)，逆为 \(\Phi_{s,t}\)。自主情形可只用经过时间 \(t-s\) 表示同一个流；一般时间依赖情形则须保留两个绝对时刻。

\begin{proposition}\label{g4:prop:0901:symplectic-flow}
	设 \((M^{2n},\omega)\) 为无边界辛流形，\(H:I\times M\to\mathbb R\) 光滑。则每个有定义的演化映射 \(\Phi_{t,s}\) 都是其开定义域到开像的辛同胚，并保持体积形式 \(\omega^n/n!\)。如果 \(M\) 紧，或存在固定紧集 \(K\subset M\) 使 \(\operatorname{supp}(\mathrm{d}_M H_t)\subset K\) 对全部 \(t\in I\) 成立，则每个 \(\Phi_{t,s}\) 对所有初值都有定义，成为 \(M\) 上的整体辛同胚。
\end{proposition}
\begin{proof}
	固定初始时刻 \(s\)。将\cref{g4:prop:0801:pullback-derivative}用于不随时间变化的形式 \(\omega\)，得到
	\[
		\frac{\mathrm{d}}{\mathrm{d}t}\Phi_{t,s}^*\omega
		=\Phi_{t,s}^*\mathcal L_{X_{H_t}}\omega
		=\Phi_{t,s}^*\mathrm{d}(\mathrm{d}_M H_t)=0.
	\]
	最后一式是逐时刻的空间外微分。因此拉回始终等于初值 \(\omega\)。逆映射与开像由演化映射的性质给出，故为辛同胚。拉回保持楔积，遂也保持 \(\omega^n/n!\)。

	若所列共同紧支条件成立，\(X_{H_t}\) 在 \(K\) 外为零，第一卷命题13.2.10保证全时间区间的演化映射存在；紧流形情形可直接取 \(K=M\)。该命题用增广场 \(\partial_t+X_{H_t}\) 的逃逸判据排除有限端点，所以并不需要假设 \(I\) 本身紧。
\end{proof}

保辛不等于能量守恒。前者在每个固定时刻由 \(\mathrm{d}_M^2H_t=0\) 保证；后者还要考虑函数本身的显式时间变化。

\begin{proposition}\label{g4:prop:0901:energy-evolution}
	设 \(H,F:I\times M\to\mathbb R\) 为无边界辛流形上的光滑函数，\(z(t)\) 满足 \(\dot z=X_{H_t}(z)\)。则
	\begin{equation}\label{g4:eq:0901:observable-evolution}
		\frac{\mathrm{d}}{\mathrm{d}t}F(t,z(t))
		=\partial_tF(t,z(t))+\mathrm{d}_M F_t(X_{H_t})(z(t)).
	\end{equation}
	特别地，
	\begin{equation}\label{g4:eq:0901:energy-evolution}
		\frac{\mathrm{d}}{\mathrm{d}t}H(t,z(t))=\partial_tH(t,z(t)).
	\end{equation}
	函数 \(F\) 沿每条局部轨迹保持不变，当且仅当 \(\partial_tF+\mathrm{d}_M F_t(X_{H_t})=0\) 在 \(I\times M\) 上恒成立。
\end{proposition}
\begin{proof}
	第一式是复合函数的链式法则。令 \(F=H\) 后，空间方向的一项为
	\[
		\mathrm{d}_M H_t(X_{H_t})=\omega(X_{H_t},X_{H_t})=0,
	\]
	得到第二式。若所列函数恒为零，第一式直接给出沿轨迹不变。反过来，每个 \((s,x)\in I\times M\) 都有一条以它为初值的局部轨迹；在初始时刻求导，便知该函数在 \((s,x)\) 处为零。
\end{proof}

沿全部轨迹保持不变的函数称为\term[diyijifen]{第一积分}[first integral]；允许显含时间时，也可用同一沿轨迹不变条件描述守恒量。对不显含时间的 \(F\)，条件化为 \(\mathrm{d}F(X_{H_t})=0\)。在标准坐标中，这一项为
\[
	\mathrm{d}F(X_{H_t})=
	\sum_i\left(
	\frac{\partial F}{\partial q^i}\frac{\partial H_t}{\partial p_i}
	-\frac{\partial F}{\partial p_i}\frac{\partial H_t}{\partial q^i}
	\right).
\]
\secref{g4:sec:0902}将把这个由两个函数产生的表达式组织成 Poisson 括号，并讨论第一积分与无穷小对称的关系。

\begin{example}\label{g4:exm:0901:time-dependent-translation}
	在 \(\mathbb R^2\) 上取 \(H_t(q,p)=a(t)p\)，其中 \(a\) 为光滑实函数。则
	\[
		X_{H_t}=a(t)\partial_q,
		\qquad
		\Phi_{t,s}(q,p)=\left(q+\int_s^t a(\tau)\,\mathrm{d}\tau,p\right).
	\]
	直接求导验证这对任意 \(s,t\in I\) 都给出全局解。它保持 \(\mathrm{d}q\wedge\mathrm{d}p\)，并保持函数 \(p\)，而沿初值动量为 \(p_0\) 的轨迹，Hamilton 函数的值是 \(a(t)p_0\)，其导数为 \(a'(t)p_0\)。若给 \(H_t\) 再加任意只依赖时间的 \(c(t)\)，速度场不变；在固定时间区间上的作用量只减去常数 \(\int_a^b c(t)\,\mathrm{d}t\)，故临界曲线不变。
\end{example}

\begin{proposition}\label{g4:prop:0901:compact-energy-complete}
	设 \((M,\omega)\) 为无边界辛流形，\(H:M\to\mathbb R\) 为光滑自主 Hamilton 函数。如果能量层 \(H^{-1}(H(x))\) 紧，则经过 \(x\) 的最大轨迹在全部实数时间上存在。特别地，若 \(H\) 为固有映射，即紧集的原像为紧集，则 \(X_H\) 完备。
\end{proposition}
\begin{proof}
	由\cref{g4:prop:0901:energy-evolution}，轨迹始终留在 \(H^{-1}(H(x))\) 中。第一卷定理13.2.2说明，无边界流形上的光滑向量场若有有限最大时间端点，其轨迹必须在接近该端点时离开每个紧集。现在整个轨迹位于同一紧能量层中，故不可能有有限左端点或右端点。若 \(H\) 固有，则每个单点 \(\{H(x)\}\) 的原像都紧，结论对所有初值成立。
\end{proof}

\begin{example}\label{g4:exm:0901:energy-with-escape}
	能量守恒本身不能保证完备。在 \(\mathbb R^2\) 上取 \(H(q,p)=q^2p\)，则
	\[
		\dot q=q^2,\qquad \dot p=-2qp.
	\]
	对初值 \(q_0>0\)，解为
	\[
		q(t)=\frac{q_0}{1-q_0t},\qquad
		p(t)=p_0(1-q_0t)^2,
		\qquad -\infty<t<1/q_0.
	\]
	代回两条方程和初值即可验证；当 \(t\uparrow1/q_0\) 时，\(q(t)\to+\infty\)，不能在 \(\mathbb R^2\) 中继续光滑延拓。另一方面，\(H(q(t),p(t))=q_0^2p_0\) 始终不变。这里能量层非紧，因而没有排除逃逸。谐振子的正能量层则为紧椭圆，前一命题也从几何上解释了其显式解为何能存在于所有时间。
\end{example}
% 纲要条目排为编号小节标题，并进入目录与书签；说明段仍保留于纲要注释中。

% !TeX root = ../../Book4.tex
% 纲要标题：### 9.2 Poisson 括号与无穷小对称
% 纲要位置：几何学计划纲要.md，第 3823 行。
% 正文状态：初稿；尚未正式审核。

% 纲要内容开始
% 
% * 括号的定义与 Jacobi 恒等式
% * 向量场括号的符号核验
% * 第一积分和约束
% * 典型动力系统
% 
% 纲要内容结束

\section{Poisson 括号与无穷小对称}
\label{g4:sec:0902}

一个 Hamilton 函数规定运动，另一个函数则可以记录运动中的某个量。沿 \(X_H\) 的轨线观察 \(f\)，其变化率是 \(\mathrm{d}f(X_H)\)。辛形式还把 \(f\) 本身变成向量场 \(X_f\)，于是这个变化率可以写成两个 Hamilton 向量场的配对。由此得到的函数运算既检验守恒量，也描述无穷小运动之间的交换关系。本节始终在无边界辛流形上讨论光滑实函数，所有流的等式都在相应解存在的时间内解释。

\subsection{括号的定义与 Jacobi 恒等式}
\label{g4:subsec:0902:01}

沿用 \(\iota_{X_f}\omega=\mathrm{d}f\)。把 \(X_g\) 代入这一等式，就得到
\[
	\mathrm{d}f(X_g)=\omega(X_f,X_g).
\]
左边是在 \(g\) 所产生的运动中观测 \(f\) 的变化率，右边则不需要选择坐标。

\begin{definition}\label{g4:def:0902:poisson-bracket}
	设 \((M,\omega)\) 为辛流形，\(f,g\in C^\infty(M,\mathbb R)\)。函数
	\begin{equation}\label{g4:eq:0902:poisson-definition}
		\{f,g\}\coloneqq\omega(X_f,X_g)=X_g(f)=-X_f(g)
	\end{equation}
	称为 \(f,g\) 的\term[poissonkuohao]{Poisson 括号}[Poisson bracket]。若 \(\{f,g\}=0\)，则称两函数\term[poissonduihe]{Poisson 对合}[in involution]。
\end{definition}

在 Darboux 坐标中，\(\omega=\sum_i\mathrm{d}q^i\wedge\mathrm{d}p_i\)，收缩方程给出
\[
	X_f=\sum_i\left(
		\frac{\partial f}{\partial p_i}\partial_{q^i}
		-\frac{\partial f}{\partial q^i}\partial_{p_i}
	\right).
\]
因此
\begin{equation}\label{g4:eq:0902:coordinate-bracket}
	\{f,g\}
	=\sum_i\left(
		\frac{\partial f}{\partial q^i}\frac{\partial g}{\partial p_i}
		-\frac{\partial f}{\partial p_i}\frac{\partial g}{\partial q^i}
	\right),
	\qquad
	\{q^i,p_j\}=\delta^i_j,\quad
	\{q^i,q^j\}=\{p_i,p_j\}=0.
\end{equation}
这也核对了括号中两个位置的次序：\(H\) 放在第二个位置时，\(\{f,H\}\) 才是沿实际 Hamilton 运动的导数。所取符号与 Cannas da Silva 的定义5.4一致。\cite[第5.1节，第140--141页]{Cannas2006Symplectic}

\begin{lemma}\label{g4:lem:0902:bracket-naturality}
	设 \(\Phi:(M,\omega_M)\to(N,\omega_N)\) 是辛同胚，\(f,g\in C^\infty(N,\mathbb R)\)。则
	\begin{equation}\label{g4:eq:0902:bracket-naturality}
		X_{f\circ\Phi}(x)
		=(\mathrm{d}\Phi_x)^{-1}X_f(\Phi(x)),
		\qquad
		\{f\circ\Phi,g\circ\Phi\}_M=\{f,g\}_N\circ\Phi.
	\end{equation}
	同一结论适用于开集之间的局部辛同胚。
\end{lemma}
\begin{proof}
	对任意 \(v\in\mathrm{T}_xM\)，链式法则和 \(\Phi^*\omega_N=\omega_M\) 给出
	\[
		\begin{aligned}
			\mathrm{d}(f\circ\Phi)_x(v)
			&=\mathrm{d}f_{\Phi(x)}(\mathrm{d}\Phi_xv)\\
			&=\omega_{N,\Phi(x)}(X_f(\Phi(x)),\mathrm{d}\Phi_xv)\\
			&=\omega_{M,x}((\mathrm{d}\Phi_x)^{-1}X_f(\Phi(x)),v).
		\end{aligned}
	\]
	非退化性确定了第一式。再对 \(f,g\) 的两个向量场应用保持配对的等式，即得第二式。
\end{proof}

\begin{theorem}\label{g4:thm:0902:poisson-identities}
	在无边界辛流形 \((M,\omega)\) 上，Poisson 括号是实双线性运算，并对任意光滑实函数 \(f,g,h\) 满足
	\begin{align}
		\{f,g\}&=-\{g,f\},\label{g4:eq:0902:antisymmetry}\\
		\{f,gh\}&=\{f,g\}h+g\{f,h\},\label{g4:eq:0902:leibniz}\\
		\{f,\{g,h\}\}+\{g,\{h,f\}\}+\{h,\{f,g\}\}&=0.
		\label{g4:eq:0902:jacobi}
	\end{align}
	最后一式称为\term[jacobihengdengshi]{Jacobi 恒等式}[Jacobi identity]。
\end{theorem}
\begin{proof}
	收缩方程线性且解唯一，所以 \(X_{af+bg}=aX_f+bX_g\)，这给出实双线性。反对称性来自 \(\omega\) 的交替性。又因为 \(\{f,gh\}=-X_f(gh)\)，对函数的乘积求导即得 Leibniz 法则\footnote{Gottfried Wilhelm von Leibniz（1646--1716），德国数学家与哲学家，研究微积分与逻辑。}；结合反对称性也得到第一个变量中的乘积法则。

	为证明 Jacobi 恒等式，先看 \(h\) 的局部流 \(\Phi_t\)。第一卷定理14.1.5的 Cartan 公式在式\eqref{g4:eq:0801:cartan}中已恢复；它给出
	\[
		\mathcal L_{X_h}\omega
		=\mathrm{d}(\iota_{X_h}\omega)+\iota_{X_h}\mathrm{d}\omega
		=\mathrm{d}^2h=0.
	\]
	由\cref{g4:prop:0801:pullback-derivative}，\(\Phi_t^*\omega=\omega\) 在流的局部定义域内成立。于是\cref{g4:lem:0902:bracket-naturality}给出
	\[
		\{f,g\}\circ\Phi_t=\{f\circ\Phi_t,g\circ\Phi_t\}.
	\]
	对 \(t\) 在零处求导。式\eqref{g4:eq:0902:coordinate-bracket}说明括号只由函数的一阶导数逐点双线性组成，故参数求导按通常乘积法则进行，得到
	\[
		X_h\{f,g\}=\{X_hf,g\}+\{f,X_hg\}.
	\]
	用 \(X_hu=\{u,h\}\) 改写，即
	\[
		\{\{f,g\},h\}=\{\{f,h\},g\}+\{f,\{g,h\}\}.
	\]
	最后使用反对称性移项，恰为式\eqref{g4:eq:0902:jacobi}。每个点附近都有上述小时间流，所以恒等式在整个 \(M\) 上成立。这里比较的是函数本身的导数，因而直接得到函数恒等式。
\end{proof}

带有上述 Lie 括号、且括号对交换乘法满足 Leibniz 法则的交换结合实代数，称为\term[poissondaishu]{Poisson 代数}[Poisson algebra]。因此 \(C^\infty(M,\mathbb R)\) 不仅能逐点相乘，还带有由辛运动确定的 Poisson 代数结构。非退化性负责把函数微分变成唯一向量场，闭性则通过 Hamilton 流保持 \(\omega\) 进入 Jacobi 恒等式的证明。

\subsection{向量场括号的符号核验}
\label{g4:subsec:0902:02}

函数的括号与向量场的 Lie 括号都描述交换，但本书的 Hamilton 约定使它们之间带有负号。直接对收缩方程求 Lie 导数，可以确定这个符号，而不必从某个例子猜测。

\begin{proposition}\label{g4:prop:0902:hamiltonian-bracket}
	设 \((M,\omega)\) 为无边界辛流形，\(f,g\in C^\infty(M,\mathbb R)\)。则
	\begin{equation}\label{g4:eq:0902:hamiltonian-bracket}
		[X_f,X_g]=-X_{\{f,g\}}.
	\end{equation}
	因此 \(f\mapsto X_f\) 为 Lie 代数反同态，而 \(f\mapsto-X_f\) 为 Lie 代数同态。
\end{proposition}
\begin{proof}
	第一卷命题13.3.11的插入恒等式为
	\[
		\iota_{[X,Y]}=\mathcal L_X\iota_Y-\iota_Y\mathcal L_X.
	\]
	其中插入固定第一个输入，Lie 导数沿正向流拉回定义。取 \(X=X_f\)、\(Y=X_g\)，并使用已证明的 \(\mathcal L_{X_f}\omega=0\)，得到
	\[
		\begin{aligned}
			\iota_{[X_f,X_g]}\omega
			&=\mathcal L_{X_f}(\mathrm{d}g)\\
			&=\mathrm{d}(X_fg)
			=-\mathrm{d}\{f,g\}
			=\iota_{-X_{\{f,g\}}}\omega.
		\end{aligned}
	\]
	第二步也可直接由 Cartan 公式作用于 \(\mathrm{d}g\) 得到。非退化性给出向量场相等。最后，
	\([{-X_f},{-X_g}]=[X_f,X_g]=-X_{\{f,g\}}\)，所以加负号后的映射为同态。
\end{proof}

例如在 \(\mathbb R^2_{q,p}\) 上取 \(f=q^2/2\)、\(g=p^2/2\)，则
\[
	X_f=-q\partial_p,\qquad X_g=p\partial_q,\qquad
	\{f,g\}=qp,
\]
而
\[
	[X_f,X_g]=-q\partial_q+p\partial_p
	=-X_{qp}.
\]
这个非零括号同时核验了函数次序和向量场交换子的次序。

\begin{proposition}\label{g4:prop:0902:central-extension}
	设 \((M,\omega)\) 为非空连通无边界辛流形。记
	\[
		\mathfrak{ham}(M,\omega)
		=\{X_f:f\in C^\infty(M,\mathbb R)\}.
	\]
	它在向量场 Lie 括号下为 Lie 代数，且有 Lie 代数正合列
	\begin{equation}\label{g4:eq:0902:central-extension}
		0\longrightarrow\mathbb R
		\xrightarrow{\;c\mapsto c\;}
		(C^\infty(M,\mathbb R),\{\ ,\ \})
		\xrightarrow{\;f\mapsto-X_f\;}
		\mathfrak{ham}(M,\omega)
		\longrightarrow0.
	\end{equation}
	左端的像包含于中间 Lie 代数的中心；这样的正合列称为\term[zhongxinkuozhang]{中心扩张}[central extension]。
\end{proposition}
\begin{proof}
	线性封闭性和括号封闭性由前一命题得到，右端映射按定义满射。又
	\[
		X_f=0\quad\Longleftrightarrow\quad \mathrm{d}f=0.
	\]
	微分为零的函数在坐标球上为常数，连通性使这些局部常数相同。因此核恰是常函数。常函数的 Hamilton 向量场为零，所以它与每个函数的 Poisson 括号都为零，证明了中心性和各处正合性。
\end{proof}

若 \(M\) 不连通，核应换成在各连通分支上分别取常值的光滑函数。即使在连通情形，保留常数也有实质作用：在 \(\mathbb R^2\) 上，\(\{q,p\}=1\)，但 \(X_q=-\partial_p\) 与 \(X_p=\partial_q\) 交换。故向量场交换只要求函数括号为常数，并不要求它为零。更具体地，把式\eqref{g4:eq:0902:central-extension}限制到 \(-X_q,-X_p\) 张成的交换 Lie 代数，这个有序基的任何一对提升都只能把 \(q,p\) 改成 \(q+a,p+b\)，它们的括号仍为 \(1\)。所以此受限扩张没有保持 Lie 括号的右逆。

\ProseParagraphBreak

第一卷定理13.3.4说明，两个向量场的括号为零，当且仅当它们的局部流在每点附近、充分小的共同时间范围内交换。因此 Poisson 对合函数产生局部交换的 Hamilton 流；在连通辛流形上，两个 Hamilton 流局部交换的精确条件是 \(\{f,g\}\) 为常数。这里不预设向量场完备，也不据此把小时间结论延伸到任意时间。

\subsection{第一积分和约束}
\label{g4:subsec:0902:03}

固定一个不随时间变化的 Hamilton 函数 \(H\)。若函数 \(f\) 沿 \(X_H\) 的每条积分曲线保持常值，按上一节的定义，\(f\) 就是该系统的第一积分\termidx[diyijifen]{第一积分}。这一条件要求对全部初值成立；某个函数在一条特殊轨线上偶然不变，还不足以成为第一积分。

\begin{proposition}\label{g4:prop:0902:first-integral}
	设 \((M,\omega)\) 为无边界辛流形，\(H,f\in C^\infty(M,\mathbb R)\)。则 \(f\) 为 \(H\) 的第一积分，当且仅当 \(\{f,H\}=0\)。在这个条件下，\(X_f\) 的局部流保持 \(H\)，并与 \(X_H\) 的局部流交换。若 \(f,g\) 都是第一积分，则其线性组合、乘积和 Poisson 括号仍是第一积分。
\end{proposition}
\begin{proof}
	沿任意积分曲线 \(\gamma\)，链式法则给出
	\[
		\frac{\mathrm{d}}{\mathrm{d}t}f(\gamma(t))
		=X_Hf(\gamma(t))=\{f,H\}(\gamma(t)).
	\]
	括号为零时函数沿曲线不变；反过来，以任意点为初值，在初始时刻求导即得该点的括号为零。这就是所述等价关系。\cite[第5.5节，命题5.12，第150--151页]{Cannas2006Symplectic}

	由反对称性，\(X_fH=\{H,f\}=0\)，所以 \(f\) 的流也保持 \(H\)。\cref{g4:prop:0902:hamiltonian-bracket}给出 \([X_f,X_H]=0\)，再由前述局部流交换定理得到交换性。线性组合与乘积的结论来自导子法则；对于括号，Jacobi 恒等式给出
	\[
		\{\{f,g\},H\}
		=\{\{f,H\},g\}+\{f,\{g,H\}\}=0.
	\]
\end{proof}

这说明第一积分同时产生保持原 Hamilton 函数的无穷小运动。特别地，守恒量的集合是 Poisson 子代数，但其中两个守恒量不必彼此对合。对于时间依赖的 \(H_t,f_t\)，相应条件则是
\[
	\partial_tf_t+\{f_t,H_t\}=0,
\]
因为沿曲线求导还包括函数自身对时间的显式依赖。

\ProseParagraphBreak

一组方程 \(F_1=c_1,\ldots,F_k=c_k\) 可以选出允许运动的状态。若这些方程在一阶上独立，就能把“运动始终满足约束”化为逐点切向条件。以下把子流形 \(S\) 称为余各向同性的，是指每点均有 \((\mathrm{T}_xS)^\omega\subset\mathrm{T}_xS\)，与第七章的线性定义一致。

\begin{proposition}\label{g4:prop:0902:regular-constraints}
	设 \((M^{2n},\omega)\) 为无边界辛流形，\(F=(F_1,\ldots,F_k):M\to\mathbb R^k\) 光滑，\(c\) 为正则值，且 \(S=F^{-1}(c)\ne\varnothing\)。则 \(S\) 为余维 \(k\) 的嵌入子流形，并有
	\begin{equation}\label{g4:eq:0902:constraint-tangents}
		\mathrm{T}_xS=\bigcap_i\ker\mathrm{d}F_i(x),
		\qquad
		(\mathrm{T}_xS)^\omega
		=\operatorname{span}\{X_{F_1}(x),\ldots,X_{F_k}(x)\}.
	\end{equation}
	此外：
	\begin{enumerate}
		\item 对光滑 Hamilton 函数 \(H\)，\(X_H\) 沿 \(S\) 切于 \(S\) 当且仅当 \(\{F_i,H\}|_S=0\) 对每个 \(i\) 成立。在这个条件下，从 \(S\) 出发的积分曲线在其整个存在时间内留在 \(S\)。
		\item \(S\) 余各向同性，当且仅当 \(\{F_i,F_j\}|_S=0\) 对所有 \(i,j\) 成立。此时 \(k\le n\)；若 \(k=n\)，则 \(S\) 为 Lagrange 子流形。
	\end{enumerate}
\end{proposition}
\begin{proof}
	第一卷定理11.4.4的正则水平集定理给出嵌入结构与第一条切空间等式：在 \(S\) 上微分 \(\mathrm{d}F\) 满射，局部坐标可把这些方程写成固定 \(k\) 个坐标。又因
	\[
		\omega(X_{F_i},v)=\mathrm{d}F_i(v)=0
		\qquad(v\in\mathrm{T}_xS),
	\]
	右端张成空间包含于辛正交补。各 \(\mathrm{d}F_i(x)\) 线性无关，非退化性说明各 \(X_{F_i}(x)\) 也线性无关。由\cref{g4:prop:0701:orthogonal-dimensions}，辛正交补的维数恰是 \(k\)，于是得到第二条等式。

	切向性要求 \(\mathrm{d}F_i(X_H)=0\)，这就是 \(\{F_i,H\}=0\) 沿 \(S\) 成立。若切向，\(X_H|_S\) 为 \(S\) 上的光滑向量场；其局部解包含到 \(M\) 后仍满足同一常微分方程。唯一性使环境解在每个经过 \(S\) 的时刻附近留在 \(S\)。而 \(S=F^{-1}(c)\) 在 \(M\) 中闭，所以一条环境积分曲线落在 \(S\) 中的时刻集合，在其连通存在区间内既开又闭。初值属于 \(S\) 时，这个集合就是整个区间。

	由式\eqref{g4:eq:0902:constraint-tangents}，余各向同性等价于每个 \(X_{F_i}(x)\) 都切于 \(S\)。再用第一项，其条件正是所有 \(\{F_j,F_i\}\) 沿 \(S\) 为零。此时辛正交补是一个 \(k\) 维各向同性子空间，第七章的维数界给出 \(k\le n\)。若 \(k=n\)，它与 \(\mathrm{T}_xS\) 同维，包含关系成为相等，故 \(S\) 为 Lagrange 子流形。
\end{proof}

只要求约束面上的括号消失就足够保证余各向同性，不必要求它们在整个 \(M\) 上消失。例如
\[
	F_1=p_1,\qquad F_2=q^1p_1+p_2
\]
在标准 \(\mathbb R^4\) 上处处微分独立，且 \(\{F_1,F_2\}=-p_1\)。共同零集恰为 \(\{p_1=p_2=0\}\)，仍是 Lagrange 子流形；两个约束函数却不在整个环境中对合。

\ProseParagraphBreak

另一方面，正则性不能省略。对标准 \(\mathbb R^2\) 上的约束 \(F=q^2=0\)，取 \(H=p\)，则 \(\{F,H\}=2q\) 在约束集合上为零，但 \(X_H=\partial_q\) 立即离开直线 \(\{q=0\}\)。这里 \(\mathrm{d}F\) 在整条直线上为零，\(\ker\mathrm{d}F\) 已经不能描述该直线的切空间。

\subsection{典型动力系统}
\label{g4:subsec:0902:04}

Poisson 括号把显式微分方程与守恒量的代数联系起来。下面两个系统分别展示：一组对合积分可以把轨线限制在较小的光滑水平集上；额外的对称性则可能产生彼此不对合的守恒量。

\begin{example}\label{g4:exm:0902:two-oscillators}
	在标准 \(\mathbb R^4_{q^1,q^2,p_1,p_2}\) 上，令
	\[
		I_j=\tfrac12\bigl((q^j)^2+p_j^2\bigr),\qquad
		H=\nu_1I_1+\nu_2I_2,\qquad \nu_1,\nu_2>0.
	\]
	两个自由度的变量彼此分离，故
	\[
		\{I_1,I_2\}=0,\qquad \{I_j,H\}=0.
	\]
	因此 \(I_1,I_2\) 均为第一积分。其微分分别为
	\(\mathrm{d}I_j=q^j\mathrm{d}q^j+p_j\mathrm{d}p_j\)，在两对坐标都不为零的开集上线性无关。对 \(c_1,c_2>0\)，共同水平集
	\[
		I_1=c_1,\qquad I_2=c_2
	\]
	是两个圆的乘积，并由\cref{g4:prop:0902:regular-constraints}成为 Lagrange 环面。Hamilton 方程逐对为
	\(\dot q^j=\nu_jp_j\)、\(\dot p_j=-\nu_jq^j\)。这些线性方程的解全时间存在，轨道保持在相应环面上。若某个 \(c_j=0\)，对应的一对坐标必须为零，该处 \(\mathrm{d}I_j=0\)，共同水平集便不能再按两个独立约束使用正则水平集结论。

	当频率相等时，还能找到不由 \(I_1,I_2\) 决定的第一积分。令
	\[
		A=q^1q^2+p_1p_2,\qquad
		B=q^1p_2-q^2p_1.
	\]
	例如 \(A\) 的 \(q\) 导数为 \((q^2,q^1)\)，\(p\) 导数为 \((p_2,p_1)\)；\(B\) 的两组导数分别为 \((p_2,-p_1)\)、\((-q^2,q^1)\)。代入坐标括号公式，得到
	\begin{equation}\label{g4:eq:0902:oscillator-brackets}
		\begin{aligned}
			\{I_1,A\}&=B,& \{I_2,A\}&=-B,\\
			\{I_1,B\}&=-A,& \{I_2,B\}&=A,\\
			\{A,B\}&=2(I_1-I_2).
		\end{aligned}
	\end{equation}
	所以
	\[
		\{A,H\}=(\nu_2-\nu_1)B,\qquad
		\{B,H\}=(\nu_1-\nu_2)A.
	\]
	当 \(\nu_1=\nu_2\) 时，\(A,B\) 都守恒，却通常不彼此对合。在固定 \(I_1,I_2>0\) 的环面上，改变两圆的相对角度会改变 \(A,B\)，所以它们记录了额外信息。它们也不是全部独立，因为直接展开给出
	\[
		A^2+B^2=4I_1I_2.
	\]
	若 \(\nu_1\ne\nu_2\)，上述 \(A,B\) 不再是整个系统的第一积分；括号公式精确记录了这个变化。
\end{example}

\begin{example}\label{g4:exm:0902:free-angular-momentum}
	取 \(\mathrm{T}^*\mathbb R^3\cong\mathbb R^6\) 上的典范辛形式和自由粒子 Hamilton 函数
	\[
		H(q,p)=\frac{|p|^2}{2m},\qquad m>0.
	\]
	其全时间流为 \((q,p)\mapsto(q+tp/m,p)\)，故三个动量分量 \(p_i\) 守恒。定义角动量
	\[
		L=q\times p,\qquad
		L_1=q^2p_3-q^3p_2,\quad
		L_2=q^3p_1-q^1p_3,\quad
		L_3=q^1p_2-q^2p_1.
	\]
	沿运动求导得到
	\[
		\dot L=(p/m)\times p=0,
	\]
	所以 \(L_i\) 也是第一积分。但这些积分之间有非零括号。若 \(\varepsilon_{ijk}\) 是由 \(\varepsilon_{123}=1\) 确定的完全反对称符号，则
	\begin{equation}\label{g4:eq:0902:angular-brackets}
		\{p_i,p_j\}=0,\qquad
		\{L_i,p_j\}=\sum_k\varepsilon_{ijk}p_k,\qquad
		\{L_i,L_j\}=\sum_k\varepsilon_{ijk}L_k.
	\end{equation}
	前两式分别来自变量独立性及
	\(\{L_i,p_j\}=\partial_{q^j}L_i\)。为核验第三式，直接计算一组：
	\[
		\begin{aligned}
			\{L_1,L_2\}
			&=(0,p_3,-p_2)\cdot(q^3,0,-q^1)
			 -(0,-q^3,q^2)\cdot(-p_3,0,p_1)\\
			&=q^1p_2-q^2p_1=L_3.
		\end{aligned}
	\]
	循环置换指标给出另外两式，再由反对称性得到全部关系。

	这些函数对应的运动也可直接辨认。对固定 \(a\in\mathbb R^3\)，令 \(L_a=a\cdot(q\times p)\)。则
	\[
		\partial_pL_a=a\times q,\qquad
		\partial_qL_a=p\times a,\qquad
		X_{L_a}(q,p)=(a\times q,a\times p).
	\]
	它同时旋转位置与动量，保持 \(|p|^2\)，因而是自由粒子系统的无穷小对称。\(\{L_a,L_b\}=L_{a\times b}\) 与\cref{g4:prop:0902:hamiltonian-bracket}相容：
	\[
		[X_{L_a},X_{L_b}]=-X_{L_{a\times b}}.
	\]
	在向量场层面也能检查这个负号：若 \(A_a(v)=a\times v\)，则线性场的 Lie 括号系数是 \(A_bA_a-A_aA_b=-A_{a\times b}\)。

	所有这些计算都在非退化的六维辛空间中进行。动量与角动量虽共有六个分量，却满足 \(p\cdot L=0\)，而且角动量分量通常彼此不对合。寻找许多守恒量、判断其微分独立性和判断其两两对合，是三个需要分别检验的问题。
\end{example}
% 纲要条目排为编号小节标题，并进入目录与书签；说明段仍保留于纲要注释中。

% !TeX root = ../../Book4.tex
% 纲要标题：### 9.3 辛同位与通量
% 纲要位置：几何学计划纲要.md，第 3830 行。
% 正文状态：初稿；尚未正式审核。

% 纲要内容开始
% 
% * 辛向量场对应闭一形式
% * Hamilton 向量场对应恰当一形式
% * 通量的路径积分及同伦变化
% * 环面平移的显式反例
% 
% 纲要内容结束

\section{辛同位与通量}
\label{g4:sec:0903}

Hamilton 方程由一个全局函数产生保辛运动。反过来，保辛运动是否总能写成 Hamilton 方程？沿运动求导后，这个问题变成一形式是否有全局原函数。局部闭性足以求解，绕闭路的积分则可能阻止原函数拼接。本节先在向量场层面识别这个障碍，再沿整条辛同位积分，得到记录运动所扫辛面积的通量。

\subsection{辛向量场对应闭一形式}
\label{g4:subsec:0903:01}

本节流形均无边界。给定光滑向量场 \(X\)，它的局部流记为 \(\phi_t\)。若所有有定义的局部流映射都保持 \(\omega\)，便称 \(X\) 为\term[xinxiangliangchang]{辛向量场}[symplectic vector field]。这里定义的是一个无穷小条件，并不预先要求场完备。

\begin{proposition}\label{g4:prop:0903:symplectic-closed}
	设 \((M,\omega)\) 为辛流形。下列三个条件等价：\(X\) 是辛向量场，\(\mathcal L_X\omega=0\)，一形式 \(\iota_X\omega\) 闭。因此收缩映射给出实向量空间同构
	\[
		\mathfrak{symp}(M,\omega)
		\longrightarrow Z^1(M),\qquad X\longmapsto\iota_X\omega,
	\]
	其中 \(\mathfrak{symp}(M,\omega)\) 表示所有辛向量场构成的空间，\(Z^1(M)=\{\alpha\in\Omega^1(M):\mathrm{d}\alpha=0\}\)。
\end{proposition}
\begin{proof}
	若局部流保辛，在 \(t=0\) 求导得到 \(\mathcal L_X\omega=0\)。反之，由\cref{g4:prop:0801:pullback-derivative}，
	\[
		\frac{\mathrm{d}}{\mathrm{d}t}\phi_t^*\omega
		=\phi_t^*\mathcal L_X\omega=0,
	\]
	所以在每个共同流域上 \(\phi_t^*\omega=\omega\)。Cartan 公式及 \(\mathrm{d}\omega=0\) 给出 \(\mathcal L_X\omega=\mathrm{d}(\iota_X\omega)\)，证明第三个条件等价。最后，\(\omega\) 非退化，所以对任意光滑一形式 \(\alpha\) 都有唯一光滑场解 \(\iota_X\omega=\alpha\)；这已在\cref{g4:lem:0801:contraction-solution}中逐点证明。限制到闭一形式便得所述同构。
\end{proof}

因而辛向量场在每个足够小的坐标球上都能写成 \(X_H\)：闭一形式在球上有原函数。不同球上的函数在连通交叠中相差常数；整体问题是能否同时调整这些常数，使局部函数成为一个全局函数。辛场与 Hamilton 场的这个区别可对照\cite[定义5.2及其后的讨论，第140页]{Cannas2006Symplectic}。

\begin{definition}\label{g4:def:0903:symplectic-isotopy}
	设 \(\phi:[0,1]\times M\to M\) 光滑，\(\phi_0=\operatorname{id}\)，每个 \(\phi_t\) 都是辛同胚。称 \((\phi_t)\) 为一条\term[xintongwei]{辛同位}[symplectic isotopy]，其速度场为
	\[
		X_t=\partial_t\phi_t\circ\phi_t^{-1}.
	\]
	这个场定义在当前所在点上，满足 \(\partial_t\phi_t(x)=X_t(\phi_t(x))\)。
\end{definition}

由拉回求导公式，辛同位的每个 \(X_t\) 都是辛向量场。反过来，若光滑时间依赖辛向量场的演化 \(\Phi_{t,s}\) 对所有初值及任意起讫时刻 \(s,t\in[0,1]\) 都有定义，则 \(\phi_t=\Phi_{t,0}\) 及其逆 \(\Phi_{0,t}\) 均全局存在，因而给出辛同位。第一卷命题13.1.9保证光滑性及反向演化；第一卷命题13.2.10说明，\(M\) 紧致或场具有共同紧支时，这个全时段存在条件成立。

\subsection{Hamilton 向量场对应恰当一形式}
\label{g4:subsec:0903:02}

在全部辛场中，Hamilton 场恰好对应于恰当一形式。于是 de Rham 上同调给出线性障碍，而无需先解动力系统。

\begin{proposition}\label{g4:prop:0903:infinitesimal-obstruction}
	设 \((M,\omega)\) 为辛流形，\(\mathfrak{ham}(M,\omega)\) 为所有 Hamilton 向量场构成的空间。序列
	\[
		0\longrightarrow\mathfrak{ham}(M,\omega)
		\longrightarrow\mathfrak{symp}(M,\omega)
		\xrightarrow{\;X\mapsto[\iota_X\omega]\;}
		H^1_{\mathrm{dR}}(M)\longrightarrow0
	\]
	正合。任意两个辛场的括号都是 Hamilton 场，具体地
	\[
		\iota_{[X,Y]}\omega=\mathrm{d}\bigl(\omega(Y,X)\bigr).
	\]
	因此上述商 Lie 代数是交换的；特别地，\(H^1_{\mathrm{dR}}(M)=0\) 时，每个辛向量场都是 Hamilton 场。
\end{proposition}
\begin{proof}
	映射的核正是 \(\iota_X\omega=\mathrm{d}H\) 的场，也就是 Hamilton 场。任取上同调类的闭代表 \(\alpha\)，由非退化性解 \(\iota_X\omega=\alpha\)，就得到其原像，所以映射满射。

	对括号，Lie 导数对张量求值的乘法法则给出
	\[
		\iota_{[X,Y]}\omega
		=\mathcal L_X(\iota_Y\omega)-\iota_Y(\mathcal L_X\omega).
	\]
	例如在向量 \(Z\) 上展开右边，两项中的 \(X(\omega(Y,Z))\) 与 \(\omega(Y,[X,Z])\) 相消，余下 \(\omega([X,Y],Z)\)。现在 \(\mathcal L_X\omega=0\)、\(\mathrm{d}\iota_Y\omega=0\)，再次使用 Cartan 公式便得到
	\[
		\iota_{[X,Y]}\omega
		=\mathrm{d}(\iota_X\iota_Y\omega)
		=\mathrm{d}\bigl(\omega(Y,X)\bigr).
	\]
	这也说明 Hamilton 场是辛场 Lie 代数的理想，且商上的括号为零。
\end{proof}

路径问题还要保证原函数随时间光滑。逐时刻存在一个函数，并不自动给出一份已经选好的光滑函数族；在一形式的情形，可以直接用路径积分完成选择。

\begin{lemma}\label{g4:lem:0903:smooth-potentials}
	设 \(M\) 连通，\(\alpha_t\) 是光滑一形式族，且每个 \(\alpha_t\) 都恰当。固定 \(x_0\in M\)，则存在唯一光滑函数 \(H:[0,1]\times M\to\mathbb R\)，使
	\[
		\mathrm{d}_xH_t=\alpha_t,\qquad H_t(x_0)=0.
	\]
\end{lemma}
\begin{proof}
	对每个 \(x\) 选一条从 \(x_0\) 到 \(x\) 的分段光滑道路，定义 \(H_t(x)=\int_{x_0}^x\alpha_t\)。因为 \(\alpha_t\) 恰当，第一卷引理14.2.5的端点积分公式说明，这个值与道路无关。为检验联合光滑性，固定一个小坐标球及其中心 \(x_1\)，先沿一条固定道路从 \(x_0\) 到 \(x_1\)，再沿球内径向线段到 \(x\)。第一段积分是 \(t\) 的光滑函数，第二段是 \((t,x)\) 的光滑参数积分。因此 \(H\) 在每个这样的球上光滑。局部径向原始公式给出 \(\mathrm{d}_xH_t=\alpha_t\)。两个满足条件的函数之差在连通 \(M\) 上为常数，并在 \(x_0\) 为零，所以唯一。
\end{proof}

\begin{definition}\label{g4:def:0903:hamiltonian-isotopy}
	若一条辛同位的速度场满足 \(\iota_{X_t}\omega=\mathrm{d}H_t\)，其中 \(H_t\) 随 \((t,x)\) 光滑，则称它为\term[Hamiltontongwei]{Hamilton 同位}[Hamiltonian isotopy]。能作为某条 Hamilton 同位的终点 \(\phi_1\) 的映射，称为\term[Hamiltontongpei]{Hamilton 同胚}[Hamiltonian diffeomorphism]，其全体记为 \(\operatorname{Ham}(M,\omega)\)。
\end{definition}

前一引理说明，在连通流形上检验给定辛同位是否为 Hamilton 同位，只须检验每个 \(\iota_{X_t}\omega\) 是否恰当。闭连通流形上还可从 \(H_t\) 减去其相对于辛体积的平均值，得到唯一满足 \(\int_M H_t\omega^n=0\) 的生成函数。该调整不改变向量场。

\begin{proposition}\label{g4:prop:0903:hamiltonian-group}
	Hamilton 同胚构成辛同胚群的正规子群。若 \(\phi_t,\psi_t\) 分别由 \(H_t,K_t\) 生成，则 \(\phi_t\circ\psi_t\) 由 \(H_t+K_t\circ\phi_t^{-1}\) 生成，\(\phi_t^{-1}\) 由 \(-H_t\circ\phi_t\) 生成。
\end{proposition}
\begin{proof}
	对固定辛同胚 \(a\)，拉回配对可得
	\[
		a_*X_H=X_{H\circ a^{-1}}.
	\]
	对复合求导，其速度为 \(X_{H_t}+(\phi_t)_*X_{K_t}\)，给出第一式。对 \(\phi_t^{-1}\phi_t=\operatorname{id}\) 求导，其速度为 \(- (\phi_t^{-1})_*X_{H_t}\)，给出第二式。两条路径都已由给定微分同胚写出，故在整个时间区间有定义。恒等映射由零函数生成，因此这些公式证明成群。最后，\(a\phi_ta^{-1}\) 由 \(H_t\circ a^{-1}\) 生成，证明正规性。
\end{proof}

\subsection{通量的路径积分及同伦变化}
\label{g4:subsec:0903:03}

以下设 \((M,\omega)\) 为闭连通辛流形。给定辛同位 \(\phi_t\)，其速度 \(X_t\) 产生闭一形式 \(\eta_t=\iota_{X_t}\omega\)。把这些无穷小障碍沿时间积分，得到一个一次上同调类。

\begin{definition}\label{g4:def:0903:flux}
	辛同位 \((\phi_t)_{0\le t\le1}\) 的\term[tongliang]{通量}[flux]定义为
	\begin{equation}\label{g4:eq:0903:flux}
		\operatorname{Flux}(\phi_t)
		=\left[\int_0^1\iota_{X_t}\omega\,\mathrm{d}t\right]
		\in H^1_{\mathrm{dR}}(M).
	\end{equation}
	积分在一形式的系数上进行；外微分与紧参数积分交换，所以所得形式闭。
\end{definition}

这里也可以使用拉回到初始点的形式 \(\alpha_t=\phi_t^*\eta_t\)。事实上，对任意固定闭一形式 \(\eta\)，随流求导给出
\begin{equation}\label{g4:eq:0903:trivial-cohomology-action}
	\phi_t^*\eta-\eta
	=\mathrm{d}\int_0^t\phi_u^*(\iota_{X_u}\eta)\,\mathrm{d}u.
\end{equation}
因此 \(\phi_t^*\) 在一次上同调上恒等。再对每个固定 \(t\) 将 \(\eta\) 取为 \(\eta_t\)，右边仍随 \(t\) 光滑，积分后可得
\[
	\operatorname{Flux}(\phi_t)
	=\left[\int_0^1\phi_t^*(\iota_{X_t}\omega)\,\mathrm{d}t\right].
\]

\begin{proposition}\label{g4:prop:0903:swept-area}
	设 \(\gamma:\mathbb R/\mathbb Z\to M\) 是光滑闭路，\(F(t,s)=\phi_t(\gamma(s))\)。给圆柱 \([0,1]\times\mathbb R/\mathbb Z\) 取定向 \(\mathrm{d}t\wedge\mathrm{d}s\)，则
	\[
		\left\langle\operatorname{Flux}(\phi_t),[\gamma]\right\rangle
		=\int_{[0,1]\times\mathbb R/\mathbb Z}F^*\omega.
	\]
	右侧是带方向的辛面积，允许圆柱的像自交或退化。
\end{proposition}
\begin{proof}
	有 \(\partial_tF=X_t(F)\)、\(\partial_sF=\mathrm{d}\phi_t(\gamma')\)，所以
	\[
		(F^*\omega)(\partial_t,\partial_s)
		=\bigl(\gamma^*\phi_t^*\iota_{X_t}\omega\bigr)(\partial_s).
	\]
	先沿 \(s\) 积分，再沿 \(t\) 积分，得到拉回版本通量在 \(\gamma\) 上的周期。它与式\eqref{g4:eq:0903:flux}代表同一个类，故结论成立。
\end{proof}

% !TeX root = ../Book4.tex
\begin{figure}[htbp]
	\centering
	\begin{tikzpicture}[font=\small,x=1.15cm,y=1.05cm]
		\fill[GeometryBlueLight] (-1.2,0) -- (-1.2,2.2) arc[start angle=180,end angle=360,x radius=1.2,y radius=0.32] -- (1.2,0) arc[start angle=0,end angle=-180,x radius=1.2,y radius=0.32];
		\draw[GeometryGray,dashed] (1.2,0) arc[start angle=0,end angle=180,x radius=1.2,y radius=0.32];
		\draw[GeometryBlue,thick,->] (-1.2,0) arc[start angle=180,end angle=355,x radius=1.2,y radius=0.32];
		\draw[GeometryBlue,thick] (0,2.2) ellipse[x radius=1.2,y radius=0.32];
		\draw[GeometryGray] (-1.2,0)--(-1.2,2.2) (1.2,0)--(1.2,2.2);
		\draw[GeometryBlue,thick,->] (0,-0.32)--(0,1.88);
		\node[right] at (0,1.1) {\(t\)};
		\node[below] at (0,-0.4) {\(s\in\mathbb R/\mathbb Z\)};
		\node[left] at (-1.2,0) {\(\gamma\)};
		\node[left] at (-1.2,2.2) {\(\phi_1\circ\gamma\)};
		\node[right,align=left] at (1.65,1.1) {\(F(t,s)=\phi_t(\gamma(s))\)\\[3pt]\(\displaystyle\int F^*\omega=\langle\operatorname{Flux},[\gamma]\rangle\)};
	\end{tikzpicture}
	\caption[闭路运动与通量]{参数圆柱的横向是闭路参数，纵向是同位时间。通量在该闭路上的取值等于按 \(\mathrm{d}t\wedge\mathrm{d}s\) 定向的拉回辛面积；图中的圆柱表示参数域，其在流形中的像可以自交。}
	\label{g4:fig:0903:swept-cylinder}
\end{figure}


\cref{g4:fig:0903:swept-cylinder}把周期计算画成一条闭路的运动：时间是圆柱的另一个方向，面积的符号由这两个方向的次序决定。保持同位两端不变而移动中间路径时，闭性将使该面积保持不变。下面直接给出形式层面的计算。

\begin{theorem}\label{g4:thm:0903:flux-homotopy}
	设 \(\phi_{s,t}\)、\(0\le s,t\le1\)，为光滑二参数辛同胚族，且 \(\phi_{s,0}=\operatorname{id}\)。令
	\[
		X_{s,t}=\partial_t\phi_{s,t}\circ\phi_{s,t}^{-1},\qquad
		Y_{s,t}=\partial_s\phi_{s,t}\circ\phi_{s,t}^{-1}.
	\]
	通量随 \(s\) 的变化由终点速度决定：
	\begin{equation}\label{g4:eq:0903:flux-variation}
		\frac{\mathrm{d}}{\mathrm{d}s}\operatorname{Flux}(\phi_{s,t})
		=[\phi_{s,1}^*(\iota_{Y_{s,1}}\omega)].
	\end{equation}
	特别地，若 \(\phi_{s,1}\) 也与 \(s\) 无关，则通量不变。
\end{theorem}
\begin{proof}
	把全部参数放入同一映射 \(G:[0,1]^2\times M\to M\)，\(G(s,t,x)=\phi_{s,t}(x)\)。令
	\[
		\alpha_{s,t}=\phi_{s,t}^*(\iota_{X_{s,t}}\omega),\qquad
		\beta_{s,t}=\phi_{s,t}^*(\iota_{Y_{s,t}}\omega),\qquad
		h_{s,t}=\omega(Y_{s,t},X_{s,t})\circ\phi_{s,t}.
	\]
	由于每张 \(\phi_{s,t}\) 保辛，完整拉回可写成
	\[
		G^*\omega=\omega+\mathrm{d}t\wedge\alpha_{s,t}
		+\mathrm{d}s\wedge\beta_{s,t}
		+h_{s,t}\,\mathrm{d}s\wedge\mathrm{d}t.
	\]
	第一卷命题14.1.8说明它闭。比较外微分中 \(\mathrm{d}s\wedge\mathrm{d}t\) 后再带一个空间输入的项，得到
	\[
		\partial_s\alpha_{s,t}-\partial_t\beta_{s,t}
		+\mathrm{d}_x h_{s,t}=0.
	\]
	因此
	\[
		\partial_s\int_0^1\alpha_{s,t}\,\mathrm{d}t
		=\beta_{s,1}-\beta_{s,0}
		-\mathrm{d}_x\int_0^1h_{s,t}\,\mathrm{d}t.
	\]
	初端恒等给出 \(Y_{s,0}=0\)，故 \(\beta_{s,0}=0\)。取上同调类就得到式\eqref{g4:eq:0903:flux-variation}。若两端都固定，进一步对 \(s\) 积分，两个通量代表之差明确是一个恰当形式。这既证明不变性，也展示了固定端点条件消去的边界项。
\end{proof}

\begin{proposition}\label{g4:prop:0903:flux-composition}
	对两条从恒等映射出发的辛同位，有
	\[
		\operatorname{Flux}(\phi_t\psi_t)
		=\operatorname{Flux}(\phi_t)+\operatorname{Flux}(\psi_t),\qquad
		\operatorname{Flux}(\phi_t^{-1})=-\operatorname{Flux}(\phi_t).
	\]
	Hamilton 同位的通量为零。任意 \(c\in H^1_{\mathrm{dR}}(M)\) 都能由某条辛同位的通量实现。
\end{proposition}
\begin{proof}
	若两速度场为 \(X_t,Y_t\)，复合的速度场为 \(X_t+(\phi_t)_*Y_t\)，且
	\[
		\iota_{X_t+(\phi_t)_*Y_t}\omega
		=\iota_{X_t}\omega+(\phi_t^{-1})^*(\iota_{Y_t}\omega).
	\]
	式\eqref{g4:eq:0903:trivial-cohomology-action}也适用于逆路径，所以取类并积分即得可加性。把 \(\psi_t\) 取为 \(\phi_t^{-1}\)，得到逆路径公式。若为 Hamilton 同位，则 \(\iota_{X_t}\omega=\mathrm{d}H_t\)，时间积分等于 \(\mathrm{d}\int_0^1H_t\,\mathrm{d}t\)，故通量为零。

	最后取 \(c\) 的闭代表 \(\alpha\)，解 \(\iota_X\omega=\alpha\)。\(M\) 紧致，故 \(X\) 完备；它的单位时间流是一条辛同位，其通量恰为 \([\alpha]=c\)。
\end{proof}

记能由辛同位连接恒等映射的辛同胚全体为 \(\operatorname{Symp}_0(M,\omega)\)。两条路径可以有同一终点而通量不同；差别来自从恒等映射出发又回到恒等映射的辛同位。定义\term[tongliangqun]{通量群}[flux group]
\[
	\Gamma_\omega
	=\{\operatorname{Flux}(\lambda_t):\lambda_0=\lambda_1=\operatorname{id}\}
	\subset H^1_{\mathrm{dR}}(M).
\]
这里路径均要求光滑；若需拼接，可先作在端点附近恒定的光滑时间重参数化，时间积分的值不变。复合和逆公式表明 \(\Gamma_\omega\) 是加法子群。对同一终点的两条路径 \(\phi_t,\psi_t\)，点态复合 \(\phi_t\psi_t^{-1}\) 是环路，其通量是两者之差。因此有良定义的群同态
\begin{equation}\label{g4:eq:0903:endpoint-flux}
	\overline{\operatorname{Flux}}:
	\operatorname{Symp}_0(M,\omega)
	\longrightarrow H^1_{\mathrm{dR}}(M)/\Gamma_\omega.
\end{equation}
该同态在 \(\operatorname{Ham}(M,\omega)\) 上为零。于是，某个终点在这个商中的通量非零，就排除了所有可能的 Hamilton 生成路径。这里使用加法群的商，定义本身不要求先证明 \(\Gamma_\omega\) 离散。

\subsection{环面平移的显式反例}
\label{g4:subsec:0903:04}

在 \(\mathbb T^2=\mathbb R^2/\mathbb Z^2\) 上取坐标模整数的 \((q,p)\)，辛形式为 \(\omega=\mathrm{d}q\wedge\mathrm{d}p\)，总面积归一化为 \(1\)。坐标本身不是全局实函数，\(\mathrm{d}q,\mathrm{d}p\) 却是全局闭一形式。沿两条基本圆周的周期分别检出它们的系数；第二卷第7.4.2小节给出
\[
	H^1_{\mathrm{dR}}(\mathbb T^2)
	=\mathbb R[\mathrm{d}q]\oplus\mathbb R[\mathrm{d}p].
\]

\begin{example}\label{g4:exm:0903:translation-path}
	取常数 \(a,b\in\mathbb R\)，考虑
	\[
		\phi_t(q,p)=(q+ta,p+tb),\qquad
		X=a\partial_q+b\partial_p.
	\]
	收缩一形式及通量为
	\[
		\iota_X\omega=a\,\mathrm{d}p-b\,\mathrm{d}q,
		\qquad \operatorname{Flux}(\phi_t)
		=a[\mathrm{d}p]-b[\mathrm{d}q].
	\]
	只要 \((a,b)\ne(0,0)\)，这个场就不是 Hamilton 场，因为至少有一个周期非零。局部函数 \(ap-bq\) 虽能生成它，却不能成为环面上的单值实函数。若 \(a,b\) 是整数，终点甚至是恒等映射：这展示了非 Hamilton 的路径仍可能有 Hamilton 终点。
\end{example}

要判断非整数平移的终点，还需排除别的路径。下面计算全部环路通量。这个计算也把上同调障碍写成轨道在覆盖空间中的平均位移。

\begin{proposition}\label{g4:prop:0903:torus-flux-lattice}
	对标准面积环面，有
	\[
		\Gamma_\omega=\mathbb Z[\mathrm{d}q]\oplus\mathbb Z[\mathrm{d}p].
	\]
	平移 \(T_{a,b}(q,p)=(q+a,p+b)\) 是 Hamilton 同胚，当且仅当 \(a,b\in\mathbb Z\)，即该平移为恒等映射。
\end{proposition}
\begin{proof}
	先记录一个平均值公式。若 \(\eta=A(q,p)\,\mathrm{d}q+B(q,p)\,\mathrm{d}p\) 闭，则 \(\partial_qB=\partial_pA\)。因此
	\[
		\frac{\mathrm{d}}{\mathrm{d}p}\int_0^1A(q,p)\,\mathrm{d}q=0,
		\qquad
		\frac{\mathrm{d}}{\mathrm{d}q}\int_0^1B(q,p)\,\mathrm{d}p=0.
	\]
	这两个周期等于 \(A,B\) 在整个单位方形的平均值。由上述一次上同调基，
	\[
		[\eta]=\left(\int_{\mathbb T^2}A\,\omega\right)[\mathrm{d}q]
		+\left(\int_{\mathbb T^2}B\,\omega\right)[\mathrm{d}p].
	\]
	设 \(\lambda_t\) 是从恒等映射出发的任意辛同位，速度写成 \(X_t=A_t\partial_q+B_t\partial_p\)。将其周期系数提升到 \(\mathbb R^2\)，解同一方程并令初值映射为恒等，得到 \(\widetilde\lambda_t\)。周期光滑场在有限时间区间有统一有界速度，所以解全时段存在；由唯一性，它与每个整数平移交换，并投影为 \(\lambda_t\)。因 \(\lambda_t\) 保持面积，换元给出
	\[
		\int_{[0,1]^2}\partial_t\widetilde\lambda_t(x)\,\mathrm{d}x
		=\int_{\mathbb T^2}X_t(\lambda_t(x))\,\omega(x)
		=\int_{\mathbb T^2}X_t(y)\,\omega(y).
	\]
	这里 \(X_t\) 的两个坐标分量都是环面上的实函数，因此最后的积分逐分量有意义。由平均值公式及 \(\iota_{X_t}\omega=A_t\mathrm{d}p-B_t\mathrm{d}q\)，若
	\[
		(D_q,D_p)=\int_{[0,1]^2}\bigl(\widetilde\lambda_1(x)-x\bigr)\,\mathrm{d}x,
	\]
	便有 \(\operatorname{Flux}(\lambda_t)=D_q[\mathrm{d}p]-D_p[\mathrm{d}q]\)。

	若 \(\lambda_1=\operatorname{id}\)，提升终点的差 \(\widetilde\lambda_1(x)-x\) 是连续的整数值向量，因而在连通 \(\mathbb R^2\) 上恒等于某个 \((k,l)\in\mathbb Z^2\)。所以每条环路的通量为 \(k[\mathrm{d}p]-l[\mathrm{d}q]\)。反过来，整数平移路径实现所有这些类，证明通量群的公式。

	若 \(\lambda_1=T_{a,b}\)，同样有 \(\widetilde\lambda_1(x)-x=(a+k,b+l)\)，其中 \(k,l\in\mathbb Z\)。若这条路径是 Hamilton 同位，其通量为零，于是 \(a+k=b+l=0\)，即 \(a,b\) 都是整数。反之，整数平移是恒等映射，由零 Hamilton 函数生成。这一证明排除了所有生成同一非整数平移的 Hamilton 路径。
\end{proof}

通量在这里同时保留了局部方程和整体周期的信息。常向量场在每张坐标图上都有线性 Hamilton 函数，但沿基本圆周的周期使局部函数无法拼接；积分到终点后，还必须把整圈运动产生的整数通量除去。这个区别将用于下一章判断一个群作用何时具有全局矩映射。
% 纲要条目排为编号小节标题，并进入目录与书签；说明段仍保留于纲要注释中。

% !TeX root = ../../Book4.tex
% 纲要标题：### 9.4* 退化 Poisson 几何
% 纲要位置：几何学计划纲要.md，第 3837 行。
% 正文状态：初稿；尚未正式审核。

% 纲要内容开始
% 
% * Poisson 张量和 Schouten 条件
% * 特征分布与辛叶
% * 线性 Lie–Poisson 结构
% * 秩跳跃与奇异叶
% 
% ---
% 
% 纲要内容结束

\OptionalSection{退化 Poisson 几何}{g4:sec:0904}

辛形式把每个协向量唯一地变成一个向量，因此每个函数都产生一个 Hamilton 向量场。如果保留函数之间的 Poisson 括号，却允许这一线性对应退化，就会出现不同维数的运动区域：有的点可以沿若干方向运动，有的点则对所有 Hamilton 系统都静止。本节从双向量场出发，先确定 Jacobi 恒等式的坐标条件，再证明这些运动区域各自带有辛形式。所有流形均为有限维、Hausdorff、第二可数的无边界光滑流形，函数均取实值。

\subsection{Poisson 张量和 Schouten 条件}
\label{g4:subsec:0904:01}
\begin{definition}\label{g4:def:0904:poisson-tensor}
	流形 \(M\) 上的光滑\term[shuangxiangliangchang]{双向量场}[bivector field]是一个截面 \(\pi\in\Gamma(\Lambda^2\mathrm{T}M)\)。它在每一点交替地配对两个协向量，并给出函数运算
	\begin{equation}\label{g4:eq:0904:bracket-tensor}
		\{f,g\}=\pi(\mathrm{d}f,\mathrm{d}g).
	\end{equation}
	若此运算对所有局部光滑函数满足 Jacobi 恒等式，则称 \(\pi\) 为\term[poissonzhangliang]{Poisson 张量}[Poisson tensor]，称 \((M,\pi)\) 为\term[poissonliuxing]{Poisson 流形}[Poisson manifold]。
\end{definition}

双线性、反对称性和两位置的 Leibniz 法则直接来自微分的乘积法则。需要检验的条件只有 Jacobi 恒等式。在坐标 \((x^1,\ldots,x^m)\) 中，写成
\begin{equation}\label{g4:eq:0904:coordinate-tensor}
	\pi=\frac12\sum_{i,j=1}^m\pi^{ij}\partial_i\wedge\partial_j,
	\qquad \pi^{ji}=-\pi^{ij},
	\qquad \{f,g\}=\sum_{i,j}\pi^{ij}(\partial_i f)(\partial_j g).
\end{equation}
因而 \(\pi^{ij}=\{x^i,x^j\}\)。系数矩阵的秩称为 \(\pi\) 在该点的秩；坐标变换只改变其矩阵表示，不改变秩。

\begin{definition}\label{g4:def:0904:sharp}
	本节以
	\begin{equation}\label{g4:eq:0904:sharp-convention}
		\alpha(\pi^\sharp\beta)=\pi(\alpha,\beta),
		\qquad \alpha,\beta\in\mathrm{T}_x^*M,
	\end{equation}
	定义丛映射 \(\pi^\sharp:\mathrm{T}^*M\to\mathrm{T}M\)。函数 \(f\) 的 Hamilton 向量场定义为 \(X_f=\pi^\sharp\mathrm{d}f\)。于是
	\begin{equation}\label{g4:eq:0904:hamiltonian-coordinate}
		X_f=\sum_{i,j}\pi^{ij}(\partial_j f)\partial_i,
		\qquad X_f(g)=\{g,f\}.
	\end{equation}
\end{definition}

这一位置次序与本章前两节一致。若 \(\omega\) 是辛形式，定义 \(\pi(\alpha,\beta)=\omega((\omega^\flat)^{-1}\alpha,(\omega^\flat)^{-1}\beta)\)，其中 \(\omega^\flat u=\iota_u\omega\)，就有 \(\pi^\sharp=(\omega^\flat)^{-1}\)。在标准坐标中，
\[
	\omega=\sum_a\mathrm{d}q^a\wedge\mathrm{d}p_a,
	\qquad \pi=\sum_a\partial_{q^a}\wedge\partial_{p_a},
	\qquad \iota_{X_f}\omega=\mathrm{d}f.
\]
\cref{g4:thm:0902:poisson-identities}已经证明由辛形式得到的 \(\pi\) 满足 Jacobi 恒等式。

\begin{proposition}\label{g4:prop:0904:jacobi-schouten}
	对坐标表示\eqref{g4:eq:0904:coordinate-tensor}，令
	\begin{equation}\label{g4:eq:0904:jacobi-coefficients}
		C^{ijk}=\sum_{\ell=1}^m\left(
			\pi^{i\ell}\partial_\ell\pi^{jk}
			+\pi^{j\ell}\partial_\ell\pi^{ki}
			+\pi^{k\ell}\partial_\ell\pi^{ij}
		\right).
	\end{equation}
	这些系数交替地依赖 \(i,j,k\)，且
	\begin{equation}\label{g4:eq:0904:jacobiator}
		\{f,\{g,h\}\}+\{g,\{h,f\}\}+\{h,\{f,g\}\}
		=\sum_{i,j,k}C^{ijk}(\partial_i f)(\partial_j g)(\partial_k h).
	\end{equation}
	特别地，\(\pi\) 是 Poisson 张量，当且仅当所有 \(C^{ijk}\) 都为零。
\end{proposition}

\begin{proof}
	将内层括号代入外层，并对乘积求导。微分落在 \(\pi^{ij}\) 上的项恰好组成右端。其余项含某个函数的二阶导数；例如含 \(\partial_i\partial_j f\) 的两组项，在利用二阶导数对称并交换哑指标后，其总和可写为
	\[
		\sum_{i,j,k,\ell}
		\bigl(\pi^{ki}\pi^{\ell j}+\pi^{\ell j}\pi^{ik}\bigr)
		(\partial_i\partial_j f)(\partial_k g)(\partial_\ell h)=0.
	\]
	含 \(g\) 或 \(h\) 的二阶导数同样抵消。这就证明\eqref{g4:eq:0904:jacobiator}。由 \(\pi^{ij}=-\pi^{ji}\) 直接交换任意两个指标，可见 \(C^{ijk}\) 交替。若所有系数为零，Jacobi 恒等式随即成立；反之取 \(f=x^i,g=x^j,h=x^k\)，就得到 \(C^{ijk}=0\)。
\end{proof}

公式\eqref{g4:eq:0904:jacobiator}还说明，其右端只依赖三个微分，而不依赖函数的其他高阶数据。左端本来就是坐标无关的运算，因此
\begin{equation}\label{g4:eq:0904:schouten-self-bracket}
	[\pi,\pi]
	=2\sum_{i<j<k}C^{ijk}\partial_i\wedge\partial_j\wedge\partial_k,
	\qquad
	\frac12[\pi,\pi](\mathrm{d}f,\mathrm{d}g,\mathrm{d}h)
	=\{f,\{g,h\}\}+\text{循环项}
\end{equation}
定义了一个全局三向量场。这是\term[schoutenkuohao]{Schouten 括号}[Schouten bracket]\footnote{Jan Arnoldus Schouten（1883--1971），荷兰数学家，研究张量分析与微分几何。}在两个相同双向量场上的取值；一般的 Schouten 括号把向量场的 Lie 括号扩展到多向量场，本节只需上述完全明确的情形。Poisson 条件由此简写为 \([\pi,\pi]=0\)。

\begin{example}\label{g4:exm:0904:jacobi-tests}
	在二维流形上，任意光滑双向量场都是 Poisson 张量，因为三向量场恒为零。局部表达式 \(a(x,y)\partial_x\wedge\partial_y\) 中的函数 \(a\) 可以有零点，也可以改变符号。三维时则已有实质条件：在 \(\mathbb R^3\) 上取
	\[
		\pi=\partial_x\wedge\partial_y+y\partial_y\wedge\partial_z.
	\]
	此时 \(\{x,y\}=1,\{y,z\}=y,\{z,x\}=0\)，所以
	\[
		\{x,\{y,z\}\}+\{y,\{z,x\}\}+\{z,\{x,y\}\}=1.
	\]
	该双向量场处处秩为 \(2\)，却不是 Poisson 张量。可见秩的条件不能代替 Jacobi 恒等式。
\end{example}

\begin{proposition}\label{g4:prop:0904:hamiltonian-invariance}
	在 Poisson 流形上，有
	\begin{equation}\label{g4:eq:0904:hamiltonian-identities}
		[X_f,X_g]=-X_{\{f,g\}},
		\qquad
		X_f\{g,h\}=\{X_f g,h\}+\{g,X_f h\}.
	\end{equation}
	每个 Hamilton 局部流保持 \(\pi\)，因而保持它的秩及 \(\operatorname{im}\pi^\sharp\)。
\end{proposition}

\begin{proof}
	在任意测试函数 \(h\) 上，用\eqref{g4:eq:0904:hamiltonian-coordinate}及 Jacobi 恒等式计算，得到
	\[
		[X_f,X_g]h
		=\{\{h,g\},f\}-\{\{h,f\},g\}
		=-\{h,\{f,g\}\}.
	\]
	第二个等式也是同一 Jacobi 恒等式的重排。由张量 Lie 导数的收缩法则，
	\[
		(\mathcal L_{X_f}\pi)(\mathrm{d}g,\mathrm{d}h)
		=X_f\{g,h\}-\{X_f g,h\}-\{g,X_f h\}=0.
	\]
	局部坐标函数的微分张成余切空间，故 \(\mathcal L_{X_f}\pi=0\)。第一卷命题13.3.7的张量拉回求导公式给出 \(\Phi_t^*\pi=\pi\)。特别地，\(\mathrm{d}\Phi_t\) 把 \(\operatorname{im}\pi_x^\sharp\) 同构地映到 \(\operatorname{im}\pi_{\Phi_t(x)}^\sharp\)，从而也保持秩。
\end{proof}

\subsection{特征分布与辛叶}
\label{g4:subsec:0904:02}
\begin{definition}\label{g4:def:0904:characteristic-distribution}
	Poisson 张量的\term[tezhengfenbu]{特征分布}[characteristic distribution]是逐点给定的子空间族
	\[
		\mathcal D_x=\operatorname{im}\pi_x^\sharp\subseteq\mathrm{T}_xM.
	\]
	它由所有 Hamilton 向量场的取值张成。其维数是 \(\pi_x\) 的秩，由交替型的标准形可知它总是偶数。
\end{definition}

这里的“分布”允许维数随点变化，因此不能直接应用第一卷的常秩 Frobenius 定理。我们先构造局部坐标，把非退化方向与剩余方向分开。证明中只用非零向量场的流盒、对易流以及逆函数定理。

\begin{theorem}\label{g4:thm:0904:local-splitting}
	设 \(\pi\) 在 \(x_0\) 处的秩为 \(2r\)。则 \(x_0\) 附近存在以该点为原点的坐标
	\[
		(q^1,\ldots,q^r,p_1,\ldots,p_r,z^1,\ldots,z^k),
		\qquad 2r+k=\dim M,
	\]
	使
	\begin{equation}\label{g4:eq:0904:local-splitting}
		\pi=\sum_{a=1}^r\partial_{q^a}\wedge\partial_{p_a}
		+\frac12\sum_{b,c=1}^k c^{bc}(z)\partial_{z^b}\wedge\partial_{z^c},
		\qquad c^{bc}(0)=0.
	\end{equation}
	横向双向量场 \(\pi_T=\frac12\sum c^{bc}(z)\partial_{z^b}\wedge\partial_{z^c}\) 本身满足 Poisson 条件。
\end{theorem}

\begin{proof}
	先设 \(r\geq1\)。可选一个局部函数 \(q\)，使 \(X_q(x_0)\neq0\)，并减去常数令 \(q(x_0)=0\)。第一卷定理13.2.7的流盒给出一个坐标函数 \(s\)，满足 \(X_qs=1\)。令 \(p=-s+s(x_0)\)，便有
	\[
		\{q,p\}=1,
		\qquad p(x_0)=0.
	\]
	这一步实际求出了与 \(q\) 配对的函数，而非只在一个切空间中选取标准基。

	记 \(A=X_p\)、\(B=-X_q\)。根据上述等式及\cref{g4:prop:0904:hamiltonian-invariance}，
	\[
		Aq=1,\quad Ap=0,\quad Bq=0,\quad Bp=1,
		\qquad [A,B]=0.
	\]
	因此 \(\mathrm{d}q,\mathrm{d}p\) 线性无关，\(T=\{q=p=0\}\) 是局部余维 \(2\) 的子流形。取其参数化 \(\sigma(z)\)，并以两个向量场的局部流定义
	\[
		F(a,b,z)=\Phi_A^a\bigl(\Phi_B^b(\sigma(z))\bigr).
	\]
	在原点处，\(F\) 的导数把 \(\partial_a,\partial_b\) 分别映为 \(A,B\)，其余方向张成 \(\mathrm{T}_{x_0}T\)，所以导数可逆。缩小定义域后，逆函数定理把 \(F\) 变为坐标图。第一卷定理13.3.4保证两个局部流对易，故这组坐标满足
	\[
		q\circ F=a,\qquad p\circ F=b,
		\qquad A=\partial_q,\qquad B=\partial_p.
	\]
	从 \(T\) 延伸而来的 \(z\) 坐标沿这两个流保持不变。

	于是 \(\{q,z^b\}=\{p,z^b\}=0\)，而 \(\{q,p\}=1\)。双向量场的坐标矩阵已分成一个标准二阶块与一个横向块。后者的系数 \(c^{bc}=\{z^b,z^c\}\) 满足
	\[
		\partial_q c^{bc}
		=X_p\{z^b,z^c\}
		=\{X_pz^b,z^c\}+\{z^b,X_pz^c\}=0,
	\]
	并且同样有 \(\partial_pc^{bc}=0\)。故它们只依赖 \(z\)。对只依赖 \(z\) 的三个函数使用 Jacobi 恒等式，就知道横向块也是 Poisson 张量。

	在原点，横向块的秩为 \(2r-2\)。对它重复刚才的构造，每次取出一个标准二阶块；有限次后横向块在原点的秩为零，得到\eqref{g4:eq:0904:local-splitting}。若最初 \(r=0\)，任意以 \(x_0\) 为原点的局部坐标已经满足结论。
\end{proof}

这个结论称为\term[weinsteinfenliedingli]{Weinstein 分裂定理}[Weinstein splitting theorem]。当 \(\pi\) 非退化时，横向部分消失；秩变化的全部可能性则留在 \(\pi_T\) 中。特别地，在以 \(x_0\) 为中心的分裂坐标内，中央子流形 \(z=0\) 上的秩恒为 \(2r\)，其切空间恰好等于 \(\mathcal D\)。

\begin{theorem}\label{g4:thm:0904:symplectic-leaves}
	每个 Poisson 流形唯一地分解为互不相交的最大连通、单射浸入的光滑流形 \(L\)，满足 \(\mathrm{T}_xL=\mathcal D_x\)。每片 \(L\) 上有唯一的辛形式 \(\omega_L\)，由
	\begin{equation}\label{g4:eq:0904:leaf-form}
		(\omega_L)_x(\pi_x^\sharp\alpha,\pi_x^\sharp\beta)
		=\pi_x(\alpha,\beta)
	\end{equation}
	给出。它们称为 \((M,\pi)\) 的\term[xinye]{辛叶}[symplectic leaf]。叶的维数等于其上 \(\pi\) 的秩，叶自身的拓扑可以比它作为 \(M\) 的子集所得的拓扑更细。
\end{theorem}

\begin{proof}
	先把逐点张成 \(\mathcal D\) 的 Hamilton 场选成一个可数族。取可数坐标开集对 \(V_i\Subset U_i\)，使 \(V_i\) 覆盖 \(M\)，且 \(\overline V_i\) 是 \(U_i\) 内的紧集。第一卷推论12.3.6给出支集包含于 \(U_i\)、在 \(\overline V_i\) 附近等于 \(1\) 的光滑截断函数 \(\chi_i\)。把 \(\chi_i x_i^j\) 在 \(U_i\) 外延拓为零，得到全局光滑函数。相应的 Hamilton 场构成一个可数族 \(Z_1,Z_2,\ldots\)，并在每一点张成 \(\mathcal D\)。

	定义 \(O_x\) 为从 \(x\) 出发、经过有限个 \(Z_j\) 的局部流所能到达的点集；每次流动只取实际存在的时间。负时间提供逆操作，有限流的连接提供传递性，故这些点集构成 \(M\) 的一个分割。所有 Hamilton 流保持 \(\pi\)，所以 \(\pi\) 的秩在 \(O_x\) 上恒为某个 \(2r\)。

	在 \(y\in O_x\) 处取\cref{g4:thm:0904:local-splitting}的坐标。其中央子流形 \(P=\{z=0\}\) 的切空间就是 \(\mathcal D\)。每个 Hamilton 场沿 \(P\) 都与 \(P\) 相切；积分曲线的唯一性保证，从 \(P\) 出发的短时间轨线留在 \(P\) 中。选取 \(2r\) 个 \(Z_j\)，使其在 \(y\) 处构成 \(\mathcal D_y\) 的一组基。依次沿它们流动，所得时间参数映射在零点的导数是这组基。逆函数定理表明其像包含 \(P\) 中 \(y\) 的一个开邻域。缩小 \(P\) 后便有 \(P\subseteq O_x\)。称这样的中央小片为一个叶片；若 \(r=0\)，叶片就是单点。

	两个叶片在交点附近可以作为同一流形的坐标图相接。确实，在交点 \(w\) 处再取分裂坐标。任一叶片在该坐标域中的点都属于 \(O_x\)，秩均为 \(2r\)，而分裂式给出
	\[
		\operatorname{rank}\pi=2r+\operatorname{rank}\pi_T(z).
	\]
	所以在这些点上 \(\pi_T(z)=0\)，\(\mathcal D\) 只有 \((q,p)\) 方向。叶片上 \(z\) 的微分为零；在含 \(w\) 的小连通片上，\(z\) 恒等于 \(0\)。两张叶片因此都局部落在同一个中央子流形内，且维数相同，故各自是其中的开集。其重叠是开集，坐标变换是环境坐标变换的限制，因而光滑。

	用全部叶片作为开坐标邻域，赋予 \(O_x\) 流形拓扑和光滑结构。其到 \(M\) 的包含映射连续且为单射浸入；环境流形的 Hausdorff 性也使两个不同的叶上点可被不相交的开集分开。还须验证第二可数性，因为单射浸入的拓扑不必等于子空间拓扑。

	有限流词 \(w=(j_1,\ldots,j_m)\) 只有可数多个。第一卷命题13.1.9给出流的开定义域及光滑性。固定这样的词，允许的时间参数因此组成 \(\mathbb R^m\) 中的开集 \(D_w\)，端点映射
	\[
		E_w(t_1,\ldots,t_m)
		=\Phi_{Z_{j_m}}^{t_m}\circ\cdots\circ\Phi_{Z_{j_1}}^{t_1}(x)
	\]
	是光滑的环境流形值映射。每个参数偏导数都是某个 \(Z_j\) 经后续 Hamilton 流推前所得的向量，故落在 \(\mathcal D\) 中。给定 \(t_0\in D_w\)，在 \(E_w(t_0)\) 处取分裂坐标，再取含 \(t_0\) 的小连通参数球，使其像在该坐标域内。所有像点的秩都是 \(2r\)，故横向块在这些点为零；于是 \(\mathrm{d}(z\circ E_w)=0\)，即整个小参数球的像落在一张中央叶片中。继续缩小参数球，可使像落在上面选取的一张叶坐标图内。

	因此 \(D_w\) 可由这样的参数球覆盖。欧氏开集有可数基，可从这一开覆盖中选出可数子覆盖；相应的可数张叶图覆盖 \(E_w(D_w)\)。再遍历可数多个流词，便得到覆盖 \(O_x\) 的可数叶图册。每张欧氏坐标图有可数基，所以 \(O_x\) 第二可数。刚才对端点映射的局部论证也说明每条流动路径在叶拓扑中连续；所有点均由有限条这类路径与 \(x\) 相连，故 \(O_x\) 连通。

	若连通单射浸入流形 \(N\) 的切空间处处等于 \(\mathcal D\)，则在任一点的分裂坐标中，同样由秩相等得到横向块为零，再由 \(\mathrm{d}z|_{\mathrm{T}N}=0\) 得到其小连通片落在中央叶片内。维数相等使这一映射成为局部微分同胚。于是 \(N\) 的各点按它们所属的 \(O_x\) 分成开集；连通性迫使只出现一片 \(O_x\)，且 \(N\) 是该叶中的开子流形。这证明最大性及唯一性。此论证也说明，可以用任何逐点张成 \(\mathcal D\) 的 Hamilton 场族来描述同一组叶。

	最后构造叶上的二形式。若 \(\pi^\sharp\alpha=0\)，则 \(\pi(\alpha,\beta)=0\) 对每个 \(\beta\) 成立，因此\eqref{g4:eq:0904:leaf-form}与协向量代表的选择无关。若 \(\pi(\alpha,\beta)=0\) 对所有 \(\beta\) 成立，则 \(\pi^\sharp\alpha=0\)，所以该二形式非退化。在每张中央叶片上，分裂式直接给出
	\[
		\omega_L=\sum_{a=1}^r\mathrm{d}q^a\wedge\mathrm{d}p_a.
	\]
	它因而光滑且闭。逐点公式又保证各叶片上的形式一致，给出所需的全局辛形式及其唯一性。
\end{proof}

在每张叶上，\eqref{g4:eq:0904:leaf-form}给出 \(\iota_{X_f|_L}\omega_L=\mathrm{d}(f|_L)\)。所以一个 Poisson 流形上的 Hamilton 系统，可以逐叶理解为通常的辛 Hamilton 系统。当 \(\pi\) 处处非退化时，叶就是 \(M\) 的连通分支，上述构造给出一个全局辛形式。这也证明了非退化 Poisson 张量与辛形式之间的双向对应。

\subsection{线性 Lie–Poisson 结构}
\label{g4:subsec:0904:03}
Lie 代数给出一类系数为线性函数的 Poisson 张量。设 \(\mathfrak g\) 是有限维实 Lie 代数。由于 \(\mathfrak g^*\) 是向量空间，\(\mathrm{T}_\xi\mathfrak g^*=\mathfrak g^*\)，从而自然有 \(\mathrm{T}_\xi^*\mathfrak g^*=\mathfrak g\)。函数在 \(\xi\) 处的微分便可以看作一个 Lie 代数元素。

\begin{proposition}\label{g4:prop:0904:lie-poisson}
	对 \(f,g\in C^\infty(\mathfrak g^*)\)，令
	\begin{equation}\label{g4:eq:0904:lie-poisson}
		\{f,g\}(\xi)=\langle\xi,[\mathrm{d}f_\xi,\mathrm{d}g_\xi]\rangle.
	\end{equation}
	它定义了 \(\mathfrak g^*\) 上的 Poisson 张量，称为\term[xianxingliepoissonjiegou]{线性 Lie--Poisson 结构}[linear Lie--Poisson structure]。若 \(\ell_a(\xi)=\langle\xi,a\rangle\)，则
	\[
		\{\ell_a,\ell_b\}=\ell_{[a,b]}.
	\]
\end{proposition}

\begin{proof}
	取 \(\mathfrak g\) 的基 \(e_1,\ldots,e_m\)，写 \([e_i,e_j]=\sum_k c_{ij}^{k}e_k\)，并以 \(x_i(\xi)=\langle\xi,e_i\rangle\) 为 \(\mathfrak g^*\) 的坐标。公式\eqref{g4:eq:0904:lie-poisson}对应于
	\[
		\pi^{ij}(x)=\sum_k c_{ij}^{k}x_k,
		\qquad
		\pi=\frac12\sum_{i,j,k}c_{ij}^{k}x_k\partial_{x_i}\wedge\partial_{x_j}.
	\]
	代入\eqref{g4:eq:0904:jacobi-coefficients}，\(C^{ijk}\) 中 \(x_s\) 的系数为
	\[
		\sum_\ell\left(
			c_{i\ell}^{s}c_{jk}^{\ell}
			+c_{j\ell}^{s}c_{ki}^{\ell}
			+c_{k\ell}^{s}c_{ij}^{\ell}
		\right).
	\]
	它正是 \([e_i,[e_j,e_k]]+[e_j,[e_k,e_i]]+[e_k,[e_i,e_j]]\) 在 \(e_s\) 方向的系数，由 Lie 代数的 Jacobi 恒等式为零。\cref{g4:prop:0904:jacobi-schouten}完成证明。
\end{proof}

若 \(\mathfrak g\) 是某个连通 Lie 群 \(G\) 的 Lie 代数，叶还具有群作用的描述。第一卷命题16.1.6--16.1.7定义了伴随表示 \(\operatorname{Ad}_g\)，并证明其无穷小作用为 \(\operatorname{ad}_a b=[a,b]\)。在对偶空间上，定义\term[yubansuizuoyong]{余伴随作用}[coadjoint action]为
\begin{equation}\label{g4:eq:0904:coadjoint-action}
	\langle\operatorname{Ad}_g^*\xi,a\rangle
	=\langle\xi,\operatorname{Ad}_{g^{-1}}a\rangle.
\end{equation}
这里的星号按公式包含了 \(g^{-1}\)，从而 \(\operatorname{Ad}_{gh}^*=\operatorname{Ad}_g^*\operatorname{Ad}_h^*\)，确实给出左作用。

\begin{proposition}\label{g4:prop:0904:coadjoint-leaves}
	\(\mathfrak g^*\) 的线性 Lie--Poisson 结构的辛叶恰好是连通群 \(G\) 的\term[yubansuiguidao]{余伴随轨道}[coadjoint orbit]。若
	\[
		a_\xi^\#=\left.\frac{\mathrm{d}}{\mathrm{d}t}\right|_{t=0}
		\operatorname{Ad}_{\exp(ta)}^*\xi,
	\]
	则轨道辛形式满足
	\begin{equation}\label{g4:eq:0904:orbit-form}
		\omega_\xi(a_\xi^\#,b_\xi^\#)=\langle\xi,[a,b]\rangle.
	\end{equation}
\end{proposition}

\begin{proof}
	对任意 \(b\in\mathfrak g\)，由伴随表示的微分得到
	\[
		a_\xi^\#(\ell_b)
		=\langle\xi,-[a,b]\rangle
		=\langle\xi,[b,a]\rangle
		=\{\ell_b,\ell_a\}(\xi).
	\]
	故 \(a^\#=X_{\ell_a}\)，它的流是 \(\xi\mapsto\operatorname{Ad}_{\exp(ta)}^*\xi\)。基元素 \(e_i\) 所给出的这些 Hamilton 场在每一点张成 \(\mathcal D\)。依照辛叶定理证明中的最后一个最大性论证，它们的有限流词可达集就是辛叶。

	另一方面，基方向的指数元生成整个连通群 \(G\)。为看清这一点，考虑
	\[
		(t_1,\ldots,t_m)\longmapsto
		\exp(t_1e_1)\cdots\exp(t_me_m).
	\]
	由第一卷定理16.2.2及命题16.2.4，指数映射在零点的导数为恒等映射，故上式的导数把标准基映到 \(e_1,\ldots,e_m\)。逆函数定理说明其像包含单位元的一个开邻域。因此由这些指数元生成的子群是开子群；其余陪集也开，连通性迫使该子群等于 \(G\)。于是有限 Hamilton 流词的可达集就是余伴随轨道。其光滑浸入结构已由辛叶定理给出，毋须预先假定该轨道在 \(\mathfrak g^*\) 中嵌入。

	最后，在叶形式\eqref{g4:eq:0904:leaf-form}中取 \(\alpha=\mathrm{d}\ell_a\)、\(\beta=\mathrm{d}\ell_b\)，即得\eqref{g4:eq:0904:orbit-form}。其良定义性、非退化性和闭性均已包含在辛叶定理中。
\end{proof}

这里自然出现的是余伴随轨道。当 Lie 代数带有适当的不变内积，可以把它与对偶空间识别，余伴随轨道才随之变为伴随轨道；对于矩阵群，后者又是矩阵共轭轨道。线性 Lie--Poisson 括号及其轨道形式也见 Cannas da Silva 的讨论。\cite[第5.1节，第141--142页]{Cannas2006Symplectic}

\begin{example}\label{g4:exm:0904:so3}
	把 \(a\in\mathbb R^3\) 对应到反对称矩阵 \(\widehat a\)，其中 \(\widehat a(v)=a\times v\)。向量三重积公式给出 \([\widehat a,\widehat b]=\widehat{a\times b}\)。内积 \(-\frac12\operatorname{tr}(AB)\) 对应于 \(\mathbb R^3\) 的点积，并在旋转共轭下不变。由 \(R\widehat aR^{-1}=\widehat{Ra}\)，\(\mathfrak{so}(3)^*\) 因而可识别为带旋转作用的 \(\mathbb R^3\)。其 Poisson 张量为
	\begin{equation}\label{g4:eq:0904:so3-tensor}
		\pi=x\partial_y\wedge\partial_z
		+y\partial_z\wedge\partial_x
		+z\partial_x\wedge\partial_y.
	\end{equation}
	直接代入括号与 \(\pi^\sharp\) 的定义可得
	\[
		\{f,g\}(\xi)=\xi\cdot(\nabla f\times\nabla g),
		\qquad X_f(\xi)=\nabla f(\xi)\times\xi.
	\]
	当 \(\xi\neq0\) 时，\(\mathcal D_\xi=\xi^\perp\)，秩为 \(2\)；原点处的秩为零。旋转在每个球面 \(S_R^2\) 上传递，故所有 \(R>0\) 的球面以及原点分别是辛叶。经上述矩阵识别，球面也就是 \(\mathfrak{so}(3)\) 中固定非零特征值的共轭轨道。

	求球面上的辛形式时，记 \(\sigma_R\) 为外法向定向的欧氏面积形式。对切向量 \(u,v\in\xi^\perp\)，有
	\begin{equation}\label{g4:eq:0904:so3-sphere-form}
		\omega_\xi(u,v)=\frac{\xi\cdot(u\times v)}{R^2}
		=\frac1R\sigma_R(u,v),
		\qquad \int_{S_R^2}\omega=4\pi R.
	\end{equation}
	确实，取 \(u=a\times\xi\)、\(v=b\times\xi\)，向量三重积给出 \(\xi\cdot(u\times v)=R^2\xi\cdot(a\times b)\)，恰与\eqref{g4:eq:0904:orbit-form}相符。因而叶的辛面积随 \(R\) 线性变化，而欧氏面积随 \(R^2\) 变化。函数 \(C(\xi)=|\xi|^2\) 满足 \(X_C=2\xi\times\xi=0\)，它在每张叶上为常数。
\end{example}

\subsection{秩跳跃与奇异叶}
\label{g4:subsec:0904:04}
分裂定理只把 \(\pi\) 在所选点处非退化的方向化为标准形式。横向块在原点为零，却可能在任意近的点重新出现非零秩。这样得到的辛叶分解，允许不同维数的叶彼此靠近。

\begin{definition}\label{g4:def:0904:casimir}
	若 \(C\in C^\infty(M)\) 满足 \(\{C,f\}=0\) 对所有局部光滑函数 \(f\) 成立，则称 \(C\) 为\term[casimirhanshu]{Casimir 函数}[Casimir function]。\footnote{Hendrik Brugt Gerhard Casimir（1909--2000），荷兰理论物理学家，研究量子力学、群表示与超导。}
\end{definition}

等价地，\(\pi^\sharp\mathrm{d}C=0\)，也就是 \(\mathrm{d}C\) 消灭整个特征分布。因此 Casimir 函数恰好是在每张连通辛叶上为常数的光滑函数。它对任意 Hamilton 系统都是第一积分。反过来，辛叶未必能由光滑 Casimir 函数的取值彼此区分。

\begin{example}\label{g4:exm:0904:rank-jump-plane}
	在 \(\mathbb R^2\) 上取
	\[
		\pi=y\partial_x\wedge\partial_y,
		\qquad
		\{f,g\}=y(f_xg_y-f_yg_x),
		\qquad
		X_f=y(f_y\partial_x-f_x\partial_y).
	\]
	由二维情形的判别，\(\pi\) 是 Poisson 张量。它在 \(y\neq0\) 处秩为 \(2\)，在横轴上秩为零。上半平面和下半平面各是一张二维辛叶，横轴上的每一个点则单独构成一张零维辛叶，见\cref{g4:fig:0904:rank-jump}。

	这个分解也可直接从运动看出。\(X_x=-y\partial_y\) 的流为 \((x,y)\mapsto(x,e^{-t}y)\)，可以在同号的非零高度之间移动；\(X_y=y\partial_x\) 的流为 \((x,y)\mapsto(x+ty,y)\)，可以在固定非零高度任意改变横坐标。因此同一个半平面内的任意两点相互可达。横轴上所有 Hamilton 场都为零，所以其上不同点不能借 Hamilton 流相互到达。

	两个开叶上的辛形式都是
	\[
		\omega=\frac1y\,\mathrm{d}x\wedge\mathrm{d}y.
	\]
	在二维开叶上它自动闭，而且直接收缩可检验 \(\iota_{X_f}\omega=\mathrm{d}f\)。其系数在接近横轴时无界，这不妨碍它在各自的叶上光滑；辛叶定理并不声称各叶的辛形式能拼成环境流形上的非退化二形式。

	若 \(C\) 是全平面上的光滑 Casimir 函数，它在两个半平面上分别为常数 \(c_+\)、\(c_-\)。在任一横轴点处使用连续性，得到该点的函数值同时等于 \(c_+\) 与 \(c_-\)。所以 \(c_+=c_-\)，且 \(C\) 在整个平面上为常数。这里虽有无穷多张零维叶，全局光滑 Casimir 函数却不能把其中任意两张分开。
\end{example}

% !TeX root = ../Book4.tex
\begin{figure}[htbp]
	\centering
	\begin{tikzpicture}[font=\small,x=1.05cm,y=1.05cm]
		\fill[GeometryBlueLight] (-3,0) rectangle (3,1.65);
		\fill[orange!12] (-3,-1.65) rectangle (3,0);
		\draw[->,GeometryGray] (-3.3,0)--(3.4,0) node[right] {\(x\)};
		\draw[->,GeometryGray] (0,-1.9)--(0,1.95) node[above] {\(y\)};
		\foreach \x in {-2.7,-2.1,-1.5,-0.9,-0.3,0.3,0.9,1.5,2.1,2.7}
			\fill[GeometryBlue] (\x,0) circle (1.7pt);
		\node[fill=white,inner sep=4pt] at (0,1.03) {\(y>0\)：一张二维叶};
		\node[fill=white,inner sep=4pt] at (0,-1.03) {\(y<0\)：一张二维叶};
		\node[fill=white,inner sep=3pt,anchor=west] at (3.55,0) {\shortstack{轴上每一点\\各是一张零维叶}};
	\end{tikzpicture}
	\caption[秩跳跃的平面 Poisson 结构]{\(\pi=y\partial_x\wedge\partial_y\) 的辛叶。两种底色各表示一个完整的二维叶；横轴上的圆点只列举其中若干零维叶，整条横轴不是一张叶。}
	\label{g4:fig:0904:rank-jump}
\end{figure}


在秩局部恒定的开集上，分裂式中的横向块处处为零，得到通常的正则叶状结构。秩变化处则须保留横向 Poisson 张量；不同秩的点集也不必各自构成光滑子流形。例如 \(\pi=xy\partial_x\wedge\partial_y\) 仍是平面 Poisson 张量，其秩零点集是两坐标轴之并，在原点不是子流形。最退化的例子是 \(\pi=0\)，此时每一点都是辛叶，每个函数都是 Casimir 函数。由此可见，退化 Poisson 几何同时记录两类信息：每张叶内部的辛运动，以及这些运动区域在整个流形中的排列方式。
% 纲要条目排为编号小节标题，并进入目录与书签；说明段仍保留于纲要注释中。


\begin{chaptersummary}
固定 \(\iota_{X_H}\omega=\mathrm{d}H\) 后，余切坐标给出 \(\dot q=H_p\)、\(\dot p=-H_q\)。作用量 \(\int(\theta(\dot z)-H)\,\mathrm{d}t\) 的变分产生相同方程，边界项则决定允许怎样移动端点。自主能量沿轨道守恒，时间依赖能量的变化由 \(\partial_tH\) 控制；能量守恒本身不排除轨道在有限时间逃向无穷远。

\ProseParagraphBreak

Poisson 括号 \(\{f,g\}=\omega(X_f,X_g)=X_gf\) 是观测量的变化率，并满足乘积法则与 Jacobi 恒等式。在本书约定下，\([X_f,X_g]=-X_{\{f,g\}}\)。第一积分构成 Poisson 子代数，却未必彼此对合；对正则约束，切向性与余各向同性都由约束函数的括号判定。

\ProseParagraphBreak

辛向量场对应闭一形式，Hamilton 向量场对应恰当一形式。通量将这份无穷小上同调类沿时间积分，也等于闭路运动扫出的带方向辛面积。它在端点固定的辛同伦下不变，并在路径复合下可加。标准二维环面的通量群是整数格，因而任何非恒等平移都不是 Hamilton 同胚，即使它局部总有 Hamilton 表达。

\ProseParagraphBreak

退化 Poisson 张量仍给每个函数一个光滑向量场，但这些场只张成某个可能变秩的分布。局部分裂把一点附近分成标准辛块与在该点为零的横向块，进而得到最大连通浸入辛叶。Lie 代数的括号在其对偶上产生线性 Poisson 结构，余伴随轨道是相应的辛叶；秩跳跃时，正维叶与点叶可以在同一空间中相邻。
\end{chaptersummary}

\BookExercises
\begin{exercise}\label{g4:ex:0901:extended-phase}
	设 \((M,\omega)\) 辛，\(I\subset\mathbb R\) 为开区间，\(H:I\times M\to\mathbb R\) 光滑。在扩充空间 \(I_t\times\mathbb R_e\times M\) 上取
	\[
		\Omega=\omega+\mathrm{d}t\wedge\mathrm{d}e,
		\qquad K(t,e,x)=e+H(t,x).
	\]
	证明 \(\Omega\) 辛，并计算自主函数 \(K\) 的 Hamilton 方程。说明它怎样恢复原时间依赖系统，以及守恒量 \(K\) 与原能量 \(H_t\) 的关系。
\end{exercise}
\begin{solution}
	\(\Omega\) 闭，且在 \(M\) 与 \((t,e)\) 两个直和方向上分别非退化，故为辛形式。令 \(X_K=a\partial_t+b\partial_e+V\)，与 \(\mathrm{d}K=\mathrm{d}e+(\partial_tH)\mathrm{d}t+\mathrm{d}_MH\) 比较收缩系数，得到
	\[
		a=1,\qquad b=-\partial_tH,\qquad V=X_{H_t}.
	\]
	若扩充系统的轨道参数为 \(\tau\)，则
	\[
		\frac{\mathrm{d}t}{\mathrm{d}\tau}=1,\qquad
		\frac{\mathrm{d}x}{\mathrm{d}\tau}=X_{H_t}(x),\qquad
		\frac{\mathrm{d}e}{\mathrm{d}\tau}=-\partial_tH(t,x).
	\]
	因此 \(t=t_0+\tau\)，投影到 \(M\) 恢复以 \(t\) 为时间的原方程。自主函数 \(K\) 守恒，也可直接计算：沿运动 \(\mathrm{d}H_t/\mathrm{d}\tau=\partial_tH\)，正好被 \(e\) 的变化抵消。任给原系统解及一个常数 \(C\)，取 \(e(t)=C-H(t,x(t))\)，便得到能量为 \(C\) 的扩充系统解。
\end{solution}

\begin{exercise}\label{g4:ex:0901:hamilton-jacobi}
	设 \(U\subset\mathbb R^n\) 为连通开集，\(I\subset\mathbb R\) 为开时间区间，\(H\in C^\infty(I\times U\times\mathbb R^n)\)。若 \(S\in C^\infty(I\times U)\) 满足\term[HamiltonJacobifangcheng]{Hamilton--Jacobi 方程}[Hamilton--Jacobi equation]
	\[
		\partial_tS(t,q)+H(t,q,\mathrm{d}_qS(t,q))=0,
	\]
	证明：沿 \(\dot q=H_p(t,q,\mathrm{d}_qS)\) 的任一局部解，取 \(p(t)=\mathrm{d}_qS(t,q(t))\)，便得到 Hamilton 轨迹。反过来，若这些随时间移动的微分图像被 Hamilton 演化局部保持，证明可从 \(S\) 减去一个仅依赖时间的函数，使方程成立；这里保持性在轨迹仍位于定义域时解释。对 \(H=|p|^2/2\)，检验 \(S=|q|^2/(2(t+c))\)，并求相应轨迹，其中 \(t+c\ne0\)。
\end{exercise}
\begin{solution}
	对方程按 \(q\) 求导，得到
	\[
		S_{tq}+H_q+S_{qq}H_p=0,
	\]
	其中 \(H\) 的各导数在 \((t,q,\mathrm{d}_qS)\) 处取值，二阶导数矩阵 \(S_{qq}\) 对称。沿题设 \(q(t)\)，
	\[
		\dot p=S_{tq}+S_{qq}\dot q=-H_q.
	\]
	与 \(\dot q=H_p\) 合在一起正是 Hamilton 方程。反过来，图像保持意味着通过图像每一点的轨道均满足 \(p=\mathrm{d}_qS\)，相同计算给出
	\[
		\mathrm{d}_q\bigl(\partial_tS+H(t,q,\mathrm{d}_qS)\bigr)=0.
	\]
	因为 \(U\) 连通，括号内是仅依赖 \(t\) 的光滑函数 \(a(t)\)。将 \(S\) 改为 \(S-\int_{t_0}^t a(u)\,\mathrm{d}u\)，其空间微分不变，且满足要求的方程。

	最后，\(\partial_tS=-|q|^2/(2(t+c)^2)\)、\(\mathrm{d}_qS=q/(t+c)\)，两项抵消。方程 \(\dot q=q/(t+c)\) 给出 \(q(t)=v(t+c)\)，\(p(t)=v\)，其中 \(v\) 为常向量。这正是一族自由粒子轨迹，公式只在不经过 \(t=-c\) 的时间区间内使用。
\end{solution}

\begin{exercise}\label{g4:ex:0901:closed-primitive-change}
	设 \(\beta\) 是 \(Q\) 上的光滑闭一形式。在 \(\mathrm{T}^*Q\) 上用 \(\vartheta=\theta+\pi^*\beta\) 代替典范一形式定义作用量。证明辛形式及固定底端点时的临界曲线不变；对闭轨迹，两种作用量之差只依赖投影闭路的光滑同伦类。取 \(Q=\mathbb R/\mathbb Z\)、\(\beta=c\,\mathrm{d}q\)，计算投影绕数为 \(k\) 时的差值。
\end{exercise}
\begin{solution}
	因 \(\mathrm{d}\beta=0\)，有 \(-\mathrm{d}\vartheta=-\mathrm{d}\theta=\omega\)。两种作用量的差为 \(\int_q\beta\)。对底曲线变分 \(q_s\)，和正文的一阶变分同样，在参数矩形上求导得
	\[
		\left.\frac{\mathrm{d}}{\mathrm{d}s}\right|_0\int_{q_s}\beta
		=[\beta(\delta q)]_a^b
		+\int_a^b\mathrm{d}\beta(\delta q,\dot q)\,\mathrm{d}t.
	\]
	固定底端点时，两项全为零，所以两种作用量具有相同的一阶变分及临界曲线。闭曲线变分中，两个端点的点与变分向量相同，边界项同样抵消；对同伦参数积分，便得到周期在光滑同伦下不变。对于圆周，把投影闭路提升为实直线中的道路，其终点减起点为整数 \(k\)，所以差值为 \(c k\)。当 \(c\ne0\) 时，闭一形式 \(\beta\) 没有全局实原函数，但仍不改变上述变分方程。
\end{solution}

\begin{exercise}\label{g4:ex:0902:cotangent-lift-generator}
	设 \(Y\) 为 \(Q\) 上的光滑向量场，定义 \(H_Y(q,p)=p(Y_q)\)。证明它的 Hamilton 流是 \(Y\) 的局部流的余切提升。再对两个向量场 \(Y,Z\) 证明
	\[
		\{H_Y,H_Z\}=-H_{[Y,Z]},\qquad
		[X_{H_Y},X_{H_Z}]=X_{H_{[Y,Z]}}.
	\]
	在 \(Q=\mathbb R\)、\(Y=q\partial_q\) 时写出显式流。
\end{exercise}
\begin{solution}
	在局部坐标中，\(H_Y=\sum_i p_iY^i(q)\)，故
	\[
		X_{H_Y}=\sum_iY^i\partial_{q^i}
		-\sum_{i,j}p_j\partial_iY^j\partial_{p_i}.
	\]
	若 \(\varphi_t\) 是 \(Y\) 的局部流，其余切提升把 \((q,p)\) 送到 \((\varphi_t(q),(\mathrm{d}\varphi_t|_q)^{-\mathsf T}p)\)。对时间求导，第一分量为 \(Y\)，第二分量由矩阵逆的求导及流的变分方程给出 \(-\mathrm{D}Y^{\mathsf T}p(t)\)，恰是所述 Hamilton 方程。初值相同，唯一性给出两个流相等，且各自在原底流有定义的范围内成立。

	由坐标括号，
	\[
		\{H_Y,H_Z\}
		=\sum_{i,j}p_j\bigl(Z^i\partial_iY^j-Y^i\partial_iZ^j\bigr)
		=-p([Y,Z]).
	\]
	再用 \([X_f,X_g]=-X_{\{f,g\}}\)，便得第二式。取 \(Y=q\partial_q\)，有 \(H_Y=qp\)，其流为 \((q,p)\mapsto(e^tq,e^{-t}p)\)。底方向的伸长与纤维方向的缩短互为倒数，保持 \(\mathrm{d}q\wedge\mathrm{d}p\)。
\end{solution}

\begin{exercise}\label{g4:ex:0901:energy-time-change}
	设 \(H\) 为辛流形上的自主 Hamilton 函数，\(f:\mathbb R\to\mathbb R\) 光滑。求 \(f\circ H\) 的向量场，并比较两系统在能量 \(H=c\) 上的轨迹、方向及周期。若 \(X_H\) 完备，证明 \(X_{f\circ H}\) 也完备。计算平面上 \(H=(q^2+p^2)/2\)、\(f(c)=c^2\) 的非静止轨道周期。
\end{exercise}
\begin{solution}
	链式法则给出 \(\mathrm{d}(f\circ H)=f'(H)\mathrm{d}H\)，由收缩方程的唯一性，
	\[
		X_{f\circ H}=f'(H)X_H.
	\]
	函数 \(H\) 沿新场也保持不变。因此从 \(H=c\) 中一点出发时，新速度就是常数 \(f'(c)\) 倍的旧速度；若旧流为 \(\phi_t\)，新流在该初值处为 \(\phi_{f'(c)t}\)。\(f'(c)>0\) 保持方向，\(f'(c)<0\) 反转方向，\(f'(c)=0\) 则整条该能量水平集都成为静止点。若旧轨道的最小正周期为 \(T\)，且 \(f'(c)\ne0\)，新周期为 \(T/|f'(c)|\)。当旧流完备时，每个初值的这个常数倍时间对所有实数 \(t\) 有定义，故新场也完备；不需要 \(f'\) 在全体能量上有统一界。

	对给定平面系统，旧非静止轨道周期为 \(2\pi\)，新速度因子为 \(2c\)。半径 \(r>0\) 的圆上 \(c=r^2/2\)，故新周期为 \(2\pi/r^2\)，方向仍为顺时针。原点静止，不赋予它最小正周期。
\end{solution}

\begin{exercise}\label{g4:ex:0902:cotangent-sphere-constraint}
	在标准 \(\mathrm{T}^*\mathbb R^{n+1}\)、\(n\ge1\)，取约束
	\[
		F_1(q,p)=\tfrac12(|q|^2-1),\qquad F_2(q,p)=q\cdot p,
		\qquad C=\{F_1=F_2=0\}.
	\]
	证明 \(C\) 是辛子流形，并以 \((q,p)\mapsto(q,p|_{\mathrm{T}_qS^n})\) 与 \(\mathrm{T}^*S^n\) 自然辛同构。对 \(H=|p|^2/2\) 限制到 \(C\) 的系统，求 Hamilton 向量场及全部轨迹；说明它为何不同于直接限制环境中的自由粒子向量场。
\end{exercise}
\begin{solution}
	在 \(C\) 上 \(|q|=1\)，两个微分的分量分别为 \((q,0)\) 和 \((p,q)\)，故线性无关，\(C\) 是余维2的正则子流形。有
	\[
		\{F_1,F_2\}=|q|^2=1\quad\text{在 }C\text{ 上}.
	\]
	由\cref{g4:prop:0902:regular-constraints}，辛正交补由 \(X_{F_1},X_{F_2}\) 张成；这两个向量的配对为 \(1\)，所以该二维补空间本身辛。于是它与自身辛正交补 \(\mathrm{T}C\) 的交为零，\(\omega|_C\) 非退化。

	给定 \(q\in S^n\)，每个 \(\mathrm{T}_qS^n\) 上的协向量都由唯一垂直于 \(q\) 的 Euclid 向量 \(p\) 表示。这个对应及其逆在球面局部标架中光滑，所以题中映射是向量丛微分同胚。两侧典范一形式对切向量的取值都为 \(p(\delta q)\)，因此映射保持典范一形式及其负外微分，证明辛同构。

	令 \(V=(p,-|p|^2q)\)。它对两个约束的导数都为零，故沿 \(C\) 切向。任一 \(W=(u,v)\in\mathrm{T}_{(q,p)}C\) 满足 \(q\cdot u=0\)，于是
	\[
		\omega(V,W)=p\cdot v+|p|^2q\cdot u=p\cdot v=\mathrm{d}(H|_C)(W).
	\]
	所以 \(V\) 就是受限 Hamilton 场。速度大小 \(r=|p|\) 沿解不变。当 \(r>0\) 时，轨迹为
	\[
		q(t)=q_0\cos(rt)+\frac{p_0}{r}\sin(rt),\qquad
		p(t)=-r q_0\sin(rt)+p_0\cos(rt).
	\]
	当 \(r=0\) 时，\((q,p)=(q_0,0)\) 静止。环境自由粒子场为 \((p,0)\)，它使 \(q\cdot p\) 的导数等于 \(|p|^2\)，在非零动量处不切于 \(C\)。受限形式所确定的场补上了保持球面约束所需的向心项。
\end{solution}

\begin{exercise}\label{g4:ex:0902:mean-zero-bracket}
	设 \((M^{2n},\omega)\) 为非空闭连通辛流形，\(n\ge1\)，记 \(\mu=\omega^n/n!\)。证明
	\[
		\int_M\{f,g\}\,\mu=0
	\]
	对任意光滑实函数成立。由此证明：若 \(X_f,X_g\) 交换，则 \(f,g\) 必须 Poisson 对合；零平均值函数的空间在 Poisson 括号下封闭，并经 \(f\mapsto-X_f\) 与 Hamilton 向量场 Lie 代数同构。比较正文在 \(\mathbb R^2\) 上的常数括号例子。
\end{exercise}
\begin{solution}
	\(M\) 紧致，故 \(X_g\) 的流 \(\phi_t\) 全时间存在并保持 \(\omega\)，从而保持 \(\mu\)。换元公式给出 \(\int_M f\circ\phi_t\,\mu=\int_M f\,\mu\)。在零时刻求导，并用紧性把导数移入积分，得到
	\[
		0=\int_M X_gf\,\mu=\int_M\{f,g\}\,\mu.
	\]
	若两向量场交换，由 \([X_f,X_g]=-X_{\{f,g\}}\) 及连通性，括号为常数。该常数的积分为零，而 \(\int_M\mu>0\)，所以常数只能是零。

	记 \(C_0^\infty(M)=\{f:\int_M f\mu=0\}\)。刚证的积分恒等式说明它对括号封闭。\(f\mapsto-X_f\) 是 Lie 代数同态，且在这个空间上的核为零。对于任意 \(X_H\)，取
	\[
		f=-H+\frac{\int_M H\mu}{\int_M\mu},
	\]
	则 \(f\) 的平均值为零，且 \(-X_f=X_H\)，证明满射。因此闭连通情形中，零平均值给出了中心扩张的一份 Lie 代数分裂。在 \(\mathbb R^2\) 上没有这里使用的有限总辛体积积分，\(\{q,p\}=1\) 与相应向量场交换可以同时成立。
\end{solution}

\begin{exercise}\label{g4:ex:0901:forward-not-onto}
	在 \((\mathbb R^2,\mathrm{d}q\wedge\mathrm{d}p)\) 上取 \(H=-q^3p\)。证明从任意初值出发的解都存在于全部 \(t\ge0\)，但每个 \(t>0\) 的时间映射都不满射到 \(\mathbb R^2\)。求该映射及其像，并直接核对它保持辛形式。说明“所有初值都能向前运行到时间1”为什么不足以产生定义在整个 \(M\) 上的辛同位。
\end{exercise}
\begin{solution}
	Hamilton 方程为 \(\dot q=-q^3\)、\(\dot p=3q^2p\)。对初值 \((q_0,p_0)\)，其解为
	\[
		\Phi_t(q_0,p_0)=\left(
		\frac{q_0}{\sqrt{1+2tq_0^2}},\;
		p_0(1+2tq_0^2)^{3/2}\right).
	\]
	直接求导可核对方程与初值，且分母对 \(t\ge0\) 始终正。若 \(t>0\)，第一分量是严格递增函数，值域为 \((-1/\sqrt{2t},1/\sqrt{2t})\)；固定第一分量后，第二分量随 \(p_0\) 遍历 \(\mathbb R\)。因此像是相应的开条带。

	记 \(A=1+2tq_0^2\)。第一分量对 \(q_0\) 的导数为 \(A^{-3/2}\)，对 \(p_0\) 的导数为零；第二分量对 \(p_0\) 的导数为 \(A^{3/2}\)。故 Jacobian 行列式为 \(1\)，即 \(\Phi_t^*(\mathrm{d}q\wedge\mathrm{d}p)=\mathrm{d}q_0\wedge\mathrm{d}p_0\)。逆映射在这个条带上为
	\[
		(q,p)\longmapsto
		\left(\frac{q}{\sqrt{1-2tq^2}},\;p(1-2tq^2)^{3/2}\right).
	\]
	它不能定义到所有目标点。辛同位要求每个时间映射都是整个 \(M\) 到自身的辛同胚，所以还需要任意目标初值的反向演化；本例只能给出到其开像的辛同胚。
\end{solution}

\begin{exercise}\label{g4:ex:0903:flux-conjugation}
	设 \((M,\omega)\) 为闭连通辛流形，\(a\) 为任意辛同胚，不要求它同位于恒等映射。对辛同位 \(\phi_t\)，证明
	\[
		\operatorname{Flux}(a\phi_ta^{-1})
		=(a^{-1})^*\operatorname{Flux}(\phi_t).
	\]
	由此证明 \((a^{-1})^*\Gamma_\omega=\Gamma_\omega\)。在标准二维环面上，对 \(A\in\operatorname{SL}(2,\mathbb Z)\) 所定义的线性辛同胚，计算它对平移速度及通量的作用。
\end{exercise}
\begin{solution}
	若 \(\phi_t\) 的速度为 \(X_t\)，共轭路径的速度为 \(a_*X_t\)。保持配对的等式给出 \(\iota_{a_*X_t}\omega=(a^{-1})^*(\iota_{X_t}\omega)\)。对时间积分并取类，就得到通量公式。共轭把基于恒等映射的辛环路送到同类环路，因此得到 \((a^{-1})^*\Gamma_\omega\subset\Gamma_\omega\)；再对 \(a^{-1}\) 使用同一论证，得到等号。

	在环面上，\(A\) 的整数系数保证映射下降，\(\det A=1\) 保持面积。若平移速度为列向量 \(v=(v_q,v_p)^{\mathsf T}\)，共轭路径直接成为速度 \(Av\) 的平移，所以通量为
	\[
		(Av)_q[\mathrm{d}p]-(Av)_p[\mathrm{d}q].
	\]
	另一方面，原通量的一形式系数向量为 \((-v_p,v_q)^{\mathsf T}\)，\((A^{-1})^*\) 把它变为 \(A^{-\mathsf T}(-v_p,v_q)^{\mathsf T}\)。二维行列式为1的等式说明这正是 \((-(Av)_p,(Av)_q)^{\mathsf T}\)，与直接计算一致。整数格被保持，也吻合通量群的不变性。
\end{solution}

\begin{exercise}\label{g4:ex:0903:commutator-hamiltonian}
	设 \((M,\omega)\) 为闭连通辛流形，\(\phi,\psi\in\operatorname{Symp}_0(M,\omega)\)。证明群交换子 \(\phi\psi\phi^{-1}\psi^{-1}\) 为 Hamilton 同胚。要求直接构造一条 Hamilton 同位，不以通量为零推断结论。
\end{exercise}
\begin{solution}
	取从恒等映射到 \(\phi\) 的辛同位 \(\phi_t\)，速度为 \(X_t\)，并令
	\[
		a_t=\phi_t\psi\phi_t^{-1},\qquad
		c_t=a_t\psi^{-1}.
	\]
	则 \(c_0=\operatorname{id}\)，\(c_1\) 是题中交换子。对复合求导，\(c_t\) 的速度为 \(V_t=X_t-(a_t)_*X_t\)，所以
	\[
		\iota_{V_t}\omega=\eta_t-(a_t^{-1})^*\eta_t,
		\qquad \eta_t=\iota_{X_t}\omega.
	\]
	\(\eta_t\) 闭。又因为 \(\psi\) 辛同位于恒等映射，固定 \(t\) 后把那条同位用 \(\phi_t\) 共轭，就给出从恒等映射到 \(a_t\) 的同位。由式\eqref{g4:eq:0903:trivial-cohomology-action}，\((a_t^{-1})^*\eta_t\) 与 \(\eta_t\) 同上同调类。因此 \(\iota_{V_t}\omega\) 在每个时刻恰当。这个一形式族光滑，\cref{g4:lem:0903:smooth-potentials}给出光滑函数族 \(K_t\)，使 \(\iota_{V_t}\omega=\mathrm{d}K_t\)。所以已构造的 \(c_t\) 是 Hamilton 同位，结论成立。
\end{solution}

\begin{optionalexercise}\label{g4:ex:0904:scaled-sphere}
	在 \(\mathbb R^3_{x,y,z}\) 上取二向量场
	\[
		\pi=z\left(
			x\,\partial_y\wedge\partial_z
			+y\,\partial_z\wedge\partial_x
			+z\,\partial_x\wedge\partial_y
		\right).
	\]
	验证它满足坐标 Jacobi 条件，并证明 \(C=x^2+y^2+z^2\) 为 Casimir 函数。确定全部辛叶；对非零半径 \(R\) 的叶，证明其辛形式在 \(\mathbf r=(x,y,z)\) 处为
	\[
		\omega_{\mathbf r}(u,v)
		=\frac{\mathbf r\cdot(u\times v)}{zR^2}.
	\]
	在球面去掉两极后的坐标 \((\theta,z)\) 中写出该形式，并证明每个二维辛叶的辛面积都无穷。
\end{optionalexercise}
\begin{solution}
	坐标括号为
	\[
		\{y,z\}=zx,\qquad \{z,x\}=zy,\qquad \{x,y\}=z^2.
	\]
	由\cref{g4:prop:0904:jacobi-schouten}的坐标判据，在三维只须核对三个不同指标的 Jacobi 式；指标重复时由反对称性自动为零。直接计算
	\[
		\begin{aligned}
			\{x,\{y,z\}\}+\{y,\{z,x\}\}+\{z,\{x,y\}\}
			&=\{x,zx\}+\{y,zy\}+\{z,z^2\}\\
			&=-xyz+xyz+0=0.
		\end{aligned}
	\]
	所以 \(\pi\) 确实为 Poisson 张量。对任意光滑函数 \(f,g\)，其括号也可写为
	\[
		\{f,g\}
		=z\,\mathbf r\cdot(\nabla f\times\nabla g).
	\]
	因为 \(\nabla C=2\mathbf r\)，有 \(\{C,g\}=0\) 对所有 \(g\) 成立，故 \(C\) 为 Casimir 函数。

	按本章的 \(\pi^\sharp\) 约定，
	\[
		X_f=\pi^\sharp(\mathrm{d}f)
		=z(\nabla f\times\mathbf r).
	\]
	在 \(z\ne0\) 处，\(\mathbf r\ne0\)，其像空间恰为 \(\mathbf r^\perp\)，维数为 \(2\)，正是球面的切空间。因此二维辛叶为
	\[
		L_R^+=\{x^2+y^2+z^2=R^2,\ z>0\},\qquad
		L_R^-=\{x^2+y^2+z^2=R^2,\ z<0\},
		\quad R>0.
	\]
	每个集合都连通，而且其切丛等于特征分布，所以是连通积分流形；任何更大的二维积分流形仍须保持 \(C\) 不变，而且不能经过特征空间为零的 \(z=0\) 平面；连通性使它留在同一个半球，故这些叶已经最大。在 \(z=0\) 平面上 \(\pi=0\)，特征空间为零，因此该平面的每一点各自构成一片零维辛叶。

	在 \(L_R^\pm\) 上，题给二形式光滑。若 \(u\ne0\) 为球面切向量，则 \(\mathbf r\times u\) 也切于球面，且
	\[
		\omega_{\mathbf r}(u,\mathbf r\times u)
		=\frac{|u|^2}{z}\ne0.
	\]
	所以形式非退化。对任意切向量 \(v\)，利用 \(\mathbf r\cdot v=0\) 及向量三重积公式，
	\[
		\begin{aligned}
			\omega_{\mathbf r}(X_f,v)
			&=\frac{\mathbf r\cdot((\nabla f\times\mathbf r)\times v)}{R^2}\\
			&=\nabla f\cdot v=\mathrm{d}(f|_{L_R^\pm})(v).
		\end{aligned}
	\]
	它与 \(\pi\) 诱导的 Hamilton 场满足所规定的收缩等式，因而正是叶上的辛形式；二维中二形式自动闭。

	令 \(a=(R^2-z^2)^{1/2}\)，球面参数化为
	\[
		\mathbf r(\theta,z)=(a\cos\theta,a\sin\theta,z),
		\qquad -R<z<R,\quad \theta\in\mathbb R/(2\pi\mathbb Z).
	\]
	计算得
	\[
		\partial_\theta\mathbf r\times\partial_z\mathbf r
		=(a\cos\theta,a\sin\theta,z)=\mathbf r,
	\]
	故
	\[
		\omega=\frac1z\,\mathrm{d}\theta\wedge\mathrm{d}z.
	\]
	两极虽不在这个坐标域内，原来的向量公式在那里仍光滑。按各叶的辛定向取正面积，位于 \(\varepsilon<|z|<R\) 的部分面积为
	\[
		2\pi\int_\varepsilon^R\frac{\mathrm{d}s}{s}
		=2\pi\log\frac R\varepsilon.
	\]
	当 \(\varepsilon\downarrow0\) 时趋于无穷。因此两个半球叶都有无穷辛面积；发散发生在趋近赤道的地方，那里恰是 \(\pi\) 退化为零的集合。
\end{solution}

\begin{optionalexercise}\label{g4:ex:0904:dense-leaves}
	设 \(\mathbb T^3=\mathbb R^3/\mathbb Z^3\)，三个坐标的周期都为 \(1\)。固定无理数 \(\alpha\)，令
	\[
		\pi=\partial_x\wedge(\partial_y+\alpha\partial_z).
	\]
	证明它是常秩 \(2\) 的 Poisson 张量，其每片辛叶均为一个单射浸入的圆柱 \(\mathbb S^1\times\mathbb R\)，在 \(\mathbb T^3\) 中稠密，却不是嵌入子流形。写出叶的参数化及辛形式，并证明所有全局光滑 Casimir 函数都是常数。
\end{optionalexercise}
\begin{solution}
	记 \(U=\partial_x\)、\(V=\partial_y+\alpha\partial_z\)。这两个全局向量场处处线性无关，故 \(\pi=U\wedge V\) 的秩处处为 \(2\)，特征分布为 \(\operatorname{span}(U,V)\)。在局部坐标中 \(\pi\) 的系数为常数，所以坐标 Jacobi 条件中的各偏导数都为零。

	过 \((x_0,y_0,z_0)\) 的叶有参数化
	\[
		\iota_{y_0,z_0}:\mathbb R/\mathbb Z\times\mathbb R\longrightarrow\mathbb T^3,
		\qquad
		\iota_{y_0,z_0}(x,t)=(x,y_0+t,z_0+\alpha t),
	\]
	右端三个分量均按模整数理解。其两个坐标导数恰为 \(U,V\)，所以这是浸入。如果两个参数有相同的像，则 \(x\) 在圆周上相同，并且
	\[
		t-t'\in\mathbb Z,\qquad \alpha(t-t')\in\mathbb Z.
	\]
	无理性迫使 \(t=t'\)，故浸入单射。在这个参数空间上，特征二向量成为 \(\partial_x\wedge\partial_t\)，相应辛形式为
	\[
		\omega_L=\mathrm{d}x\wedge\mathrm{d}t.
	\]
	这里 \(\mathrm{d}x\) 是周期 \(1\) 圆周上的全局一形式。任一切于特征分布的路径，在局部提升的坐标中都满足 \(\mathrm{d}z=\alpha\,\mathrm{d}y\)；沿参数区间拼接这些提升后，它始终落在上述同一像中。因此这个像是最大连通积分流形，其叶拓扑由所给圆柱参数化确定。

	下面证明稠密性。先证明 \(\{n\alpha\bmod1:n\in\mathbb Z\}\) 在圆周上稠密。对每个正整数 \(N\)，把 \([0,1)\) 分成 \(N\) 个等长区间。数 \(0,\alpha,\ldots,N\alpha\) 的小数部分中，至少两个落在同一区间。取其差，必要时反号，就得到整数 \(k\ne0,m\)，使
	\[
		0<\delta=k\alpha-m<\frac1N.
	\]
	严格正性由 \(\alpha\) 无理保证。对任意 \(\beta\in[0,1)\)，令 \(j=\lfloor\beta/\delta\rfloor\)，则
	\[
		0\le\beta-j\delta<\delta,\qquad 0\le j\delta<1,
	\]
	而 \(j\delta\) 正是 \(jk\alpha\) 的小数部分。让 \(N\) 增大，就能用整数倍的 \(\alpha\) 任意逼近 \(\beta\)，证明了圆周稠密性。

	给定目标点的两个后坐标 \((y_*,z_*)\)，取 \(t=y_*-y_0+n\)，其中 \(n\in\mathbb Z\)。第二坐标恰为 \(y_*\)，第三坐标为
	\[
		z_0+\alpha(y_*-y_0)+n\alpha\pmod1,
	\]
	因而能任意逼近 \(z_*\)。第一坐标 \(x\) 可任取，所以整片叶在三维环面中稠密。

	它不是嵌入。由刚才的逼近，可选非零整数 \(n_j\)，使 \(n_j\alpha\bmod1\to0\)，并使 \(|n_j|\to\infty\)：任何有限组非零整数倍都与整数保持正距离，无理性保证这一距离不为零。因此
	\[
		\iota_{y_0,z_0}(x_0,n_j)\longrightarrow
		\iota_{y_0,z_0}(x_0,0)
	\]
	在环境环面中成立，而 \((x_0,n_j)\) 不在圆柱中趋于 \((x_0,0)\)。到其像的逆映射不连续，故该单射浸入不是嵌入。

	最后，若 \(C\) 为全局光滑 Casimir 函数，则
	\[
		0=\pi^\sharp(\mathrm{d}C)=V(C)U-U(C)V.
	\]
	线性无关性给出 \(U(C)=V(C)=0\)，所以 \(C\) 在每片连通叶上为常数。任取一片稠密叶，由 \(C\) 连续，该常数便延伸到整个 \(\mathbb T^3\)。反之，常函数当然都是 Casimir 函数。
\end{solution}

\begin{optionalexercise}\label{g4:ex:0904:nilpotent-linear}
	设三维实 Lie 代数有基 \(e_1,e_2,e_3\)，唯一非零的基本括号为 \([e_1,e_2]=e_3\)。在对偶空间上取对应线性坐标 \(x,y,z\)，采用 \(\{\ell_u,\ell_v\}=\ell_{[u,v]}\) 的约定。
	\begin{enumerate}
		\item 写出线性 Poisson 张量，确定所有辛叶及其辛形式。
		\item 证明全局光滑 Casimir 函数恰为 \(F(z)\)，其中 \(F\in C^\infty(\mathbb R,\mathbb R)\)。
		\item 对 \(H=(x^2+y^2)/2\)，求每条平面 \(z=c\) 上的 Hamilton 轨道，说明非平衡轨道的周期、旋转方向以及 \(c=0\) 时的情形。
	\end{enumerate}
\end{optionalexercise}
\begin{solution}
	1. 基本括号给出 \(\{x,y\}=z\)、\(\{y,z\}=\{z,x\}=0\)，所以
	\[
		\pi=z\,\partial_x\wedge\partial_y.
	\]
	这也直接满足坐标 Jacobi 条件，因为唯一需检验的坐标 Jacobi 式为 \(\{z,z\}=0\)。当 \(z=c\ne0\) 时，特征分布由 \(\partial_x,\partial_y\) 张成，因而整张平面 \(z=c\) 是一片辛叶；其辛形式为
	\[
		\omega_c=\frac1c\,\mathrm{d}x\wedge\mathrm{d}y.
	\]
	在 \(z=0\) 平面上，张量为零，每一点各自是一片零维辛叶。因此整张 \(z=0\) 平面并不是一片二维叶。

	2. 若 \(C\) 是光滑 Casimir 函数，则
	\[
		\{C,x\}=-z\,\partial_yC=0,\qquad
		\{C,y\}=z\,\partial_xC=0.
	\]
	在每张 \(z=c\ne0\) 的平面上，\(C\) 的两个切向偏导数均为零，所以 \(C\) 在该连通平面上为常数。令 \(F(z)=C(0,0,z)\)，它是光滑实函数，且在 \(z\ne0\) 处已有 \(C(x,y,z)=F(z)\)。固定任意 \(x,y\) 并令 \(z\to0\)，连续性给出 \(C(x,y,0)=F(0)\)，故等式在全空间成立。反过来，由 \(\pi^\sharp(\mathrm{d}z)=0\)，每个 \(F(z)\) 都是 Casimir 函数。虽然 \(z=0\) 上有许多不同的点叶，全局光滑 Casimir 函数在这些点上的值仍必须相同。

	3. 按 \(\iota_{X_H}\omega_c=\mathrm{d}(H|_{z=c})\)，或者直接用 \(\pi^\sharp\mathrm{d}H\)，得到
	\[
		X_H=zy\,\partial_x-zx\,\partial_y,\qquad
		\dot x=cy,\quad \dot y=-cx,\quad \dot z=0.
	\]
	因此以 \((x_0,y_0,c)\) 为初值的轨道为
	\[
		\begin{aligned}
			x(t)&=x_0\cos(ct)+y_0\sin(ct),\\
			y(t)&=y_0\cos(ct)-x_0\sin(ct),\\
			z(t)&=c.
		\end{aligned}
	\]
	若 \(c\ne0\) 且 \(x_0^2+y_0^2>0\)，轨道是圆，最小正周期为 \(2\pi/|c|\)。相对于 \(x\) 向右、\(y\) 向上的通常平面定向，\(c>0\) 时顺时针，\(c<0\) 时逆时针；例如在正 \(x\) 轴上，速度的 \(y\) 分量为 \(-cx\)，直接确定了方向。每张非零平面上的原点 \((0,0,c)\) 是平衡点。当 \(c=0\) 时，整个张量以及 \(X_H\) 都为零，全部轨道静止。
\end{solution}
